% Theorem A: bar-cobar reconstruction, ordered descent, and Verdier duality.

\providecommand{\Bbarch}{\bar{B}^{\mathrm{ch}}}
\providecommand{\Omegach}{\Omega^{\mathrm{ch}}}
\providecommand{\Kosz}{\mathsf{Kosz}}
\providecommand{\Conil}{\mathsf{Conil}}
\providecommand{\SpCh}{\mathrm{Sp}^{\mathrm{ch}}}
\providecommand{\PhiFA}{\Phi^{\mathrm{FA}}}
\providecommand{\FactOp}{\mathsf{FactOp}}
\providecommand{\FactProp}{\mathsf{FactProp}}
\providecommand{\CoFact}{\mathsf{CoFact}}
\providecommand{\CoFactProp}{\mathsf{CoFactProp}}
\providecommand{\Dmod}{D\text{-}\mathsf{mod}}
\providecommand{\Tc}{\mathrm{T}^{\mathrm c}}
\providecommand{\rel}{\mathrm{rel}}
\providecommand{\PTVV}{\mathrm{PTVV}}
\providecommand{\CY}{\mathrm{CY}}
\providecommand{\Arnold}{\mathrm{Arnold}}

\chapter[Theorem $A$: Bar--Cobar Reconstruction and Verdier Duality]%
{Theorem $A$: Bar--Cobar Reconstruction and Verdier Duality}
\label{ch:theorem-A-infinity-2}\label{chap:theorem-A-infinity-2}

\section*{Notation, hypotheses, and lane structure}

The full bar construction and the quadratic Koszul dual occupy two
successive layers.  For an augmented dg algebra $A$, the full
tensor coalgebra $B(A)=T^c(s^{-1}\bar A)$ gives a functorial
resolution
\[
 \varepsilon_A\colon\Omega B(A)\xrightarrow{\ \sim\ }A
\]
for every $A$ \cite[Corollary~2.3.4]{LV12}.  Under the quadratic
Koszul package, the canonical comparison
$A^i\rightarrow B(A)$ is a quasi-isomorphism; finite or continuous
duality then produces $A^!$.  Higher Massey operations occur as
components of the full bar differential and measure the passage from
the quadratic model to the full resolution.

On the Ran space there is a second distinction.  Francis--Gaitsgory's
Proposition~4.1.2 gives bar--cobar equivalence for a reduced operad in
any pro-nilpotent stable symmetric-monoidal category, while the
pro-nilpotence argument in the proof of their Theorem~5.1.1 applies
this result to the chiral category $D(\Ran X)$
\cite{Francis2012}.
Their factorization category is the full subcategory of chiral
commutative coalgebras whose structure maps are equivalences
\cite[Definition~2.4.7]{Francis2012}.  The published associative
bar--cobar theorem lives in $D(\Ran X)$.  Its restriction to the
associative factorization objects of this volume is the closure theorem
named $H_{\mathrm{fact}}$ below.

The notation is fixed accordingly:
\[
 R_X(A):=\Omega_XB_X(A),
 \qquad
 K_X(A):=\mathbb D_{\Ran}B_X(A).
\]
$R_X$ denotes reconstruction by the full bar.  $K_X$ denotes
Verdier-dualized Koszul duality on the stated dualizable domain.  For
a connected positive quadratic presentation $A=T_X(V)/(R)$, the five
typed objects are
\[
 A,\quad B_X(A),\quad
 A^{\mathrm i}=C_X(s^{-1}V,s^{-2}R),\quad
 K_X(A)=\mathbb D_{\Ran}B_X(A),\quad
 Z_{\ch}^{\mathrm{der}}(A)=C^\bullet_{\ch}(A,A).
\]
The quadratic presentation carries its canonical twisting morphism
\[
 \tau_{\mathrm i}\colon A^{\mathrm i}\longrightarrow A.
\]
Their comparison maps have distinct targets:
\[
 q_A\colon A^{\mathrm i}\longrightarrow B_X(A),\qquad
 \mathbb D(q_A)\colon K_X(A)\longrightarrow
 \mathbb D_{\Ran}(A^{\mathrm i}),\qquad
 \nu_A\colon K_X(A)\longrightarrow A^!.
\]
The first is quadratic recognition, the second is its Verdier-dual
map, and the third is the chosen comparison with the strict quadratic
partner.  This chapter
proves the ambient reconstruction, isolates the factorization closure
theorem, and states the quadratic and Verdier refinements on their
typed domains.  The categorical level is $(\infty,1)$; the
factorization-enriched properadic transfer is governed by
$H_{\mathrm{prop}}$.

\paragraph{Ran ambient and four Theorem~A packages.}
Let $X$ be a separated finite-type scheme over a characteristic-zero
field.  Francis--Gaitsgory's ambient datum is
\[
 \mathcal C_X=\bigl(D(\Ran X),\otimes^{\ch}\bigr),
 \qquad \mathcal O=\mathsf{Ass}_{\mathrm{red}}.
\]
The proof of \cite[Theorem~5.1.1]{Francis2012} makes
$\mathcal C_X$ pro-nilpotent, and
\cite[Proposition~4.1.2]{Francis2012} then supplies enhanced
bar--cobar equivalence for $\mathcal O$, with divided-power
coalgebras on the coalgebra side.  These are the published ambient
data.  Four packages govern the subsequent Theorem~A transitions used in this
volume.
\begin{itemize}
\item[$H_{\mathrm{fact}}$] \emph{Factorization-restriction package.}
The chosen associative factorization subcategories satisfy four
closure conditions: enhanced bar and cobar preserve the
factorization equivalences on every disjoint configuration locus;
their collision maps satisfy base change along the diagonal
stratification; units and completed inverse limits remain in the
chosen subcategories; and a regular-singular ordered local system has
Fulton--MacPherson nearby cycles compatible with the chiral collision
product.  The last clause is denoted $H_{\mathrm{fact}}^{R}$ when an
$R$-matrix is present.  These closure conditions form a conditional
input of this volume.
\item[$H_{\mathrm{VD}}$] \emph{Verdier package.}
$B_X(A)$ lies in a constructible or holonomic dualizable subcategory
on which $\mathbb D_{\Ran}$ is defined and exchanges the relevant
coalgebra and algebra structures.  Every base-change and K\"unneth
map used after duality belongs to this package.  The properadic lane
also includes a specified evaluation--coevaluation pair between $A$
and $K_X(A)$ satisfying the two zig--zag identities.
\item[$H_{\mathrm{conv}}$] \emph{Convergence package.}
The increasing bar-length filtration on the direct-sum bar is
exhaustive; a completed algebra satisfies
$A\simeq\varprojlim_q A/\bar A^{\,q}$; and a completed coalgebra has a
pro-coradical presentation
$C\simeq\varprojlim_q C_{\leq q}$ whose transition maps are
cohomologically Mittag--Leffler in every input--output profile.
These are three independent clauses, denoted respectively
$H_{\mathrm{len}}$, $H_{\mathrm{aug}}$, and $H_{\mathrm{ML}}$.
\item[$H_{\mathrm{prop}}$] \emph{Factorization-properad transfer
package.}
The factorization enrichment satisfies the monoid axiom, has
cofibrant unit, admits connected-graph substitution, carries weak
equivalences through the transfer functors, and preserves
factorization descent.  The full formulation appears in
Theorem~\ref{thm:factorization-properad-transfer}.
\end{itemize}
The quadratic comparison is the Theorem~B interface and carries
$H_{\mathrm{CL}}(A,A^{\mathrm i},\tau_{\mathrm i})$ from
Definition~\ref{def:chiral-comparison-lemma-package}.  Finite-weight
models provide concrete instances of $H_{\mathrm{VD}}$ and
$H_{\mathrm{conv}}$.
The universal full bar--cobar counit belongs to the published ambient
theorem.  Factorization closure, completed chain realization,
quadratic compression, Verdier transport, and properadic transfer
carry the five displayed packages.  Weak curvature retains
Positselski's CDG definition as a separate structure.

\paragraph{Parametric strength taxonomy.}
\label{par:A-strength-taxonomy}%
\index{Theorem A!parametric strength}%
Theorem~A has three typed layers.
\begin{itemize}
\item[\textsc{Ran reconstruction; $\ClaimStatusProvedElsewhere$.}]
\emph{Type signature:} (Open quadrant; chiral Ran presentation;
Beilinson level $1\leftrightarrow2$; $(\infty,1)$; pro-nilpotent
ambient $\mathcal C_X$; reduced associative operad).  Enhanced bar and cobar form mutually inverse
equivalences for the reduced associative operad in
$D(\Ran X)$ \cite[Proposition~4.1.2]{Francis2012}.
\item[\textsc{Factorization closure; $\ClaimStatusConjectured$.}]
\emph{Type signature:} (Open quadrant; associative factorization
presentation; Beilinson level $1\leftrightarrow2$; pro-nilpotent
ambient $\mathcal C_X$; package $H_{\mathrm{fact}}$).  The closure theorem
restricts enhanced bar and cobar to the factorization subcategories
used in this volume, including collision and completion data.
\item[\textsc{Theorem B interface and Verdier refinement;
$\ClaimStatusConditional$.}]
\emph{Type signature:} (Open quadrant; quadratic and Verdier
presentations; Beilinson level $2$; packages
$H_{\mathrm{CL}}(A,A^{\mathrm i},\tau_{\mathrm i})\cup H_{\mathrm{VD}}$).
The comparison
$A^i\to B_X(A)$ and Verdier duality identify the full bar model with
the quadratic and dual models on this locus.
\end{itemize}
Thus $R_X=\Omega_XB_X$ is the reconstruction functor.  The functor
$K_X=\mathbb D_{\Ran}B_X$ carries the Verdier--Koszul refinement;
biduality supplies its involutive form on the dualizable locus.

\paragraph{Lane split: Priddy versus Positselski.}
\label{par:A-priddy-vs-positselski}%
\index{Theorem A!Priddy lane vs Positselski lane}%
\index{Priddy lane!chain-level Theorem A}%
\index{Positselski lane!coderived Theorem A}%
Three constructions meet along explicit comparison morphisms.
\begin{itemize}
\item[\textsc{Full bar lane.}]
The tensor coalgebra $B(A)$ and the counit
$\Omega B(A)\xrightarrow{\sim}A$ give the functorial resolution
\cite[Corollary~2.3.4]{LV12}.  Francis--Gaitsgory supply its enhanced
form in the pro-nilpotent chiral Ran ambient.
\item[\textsc{Quadratic lane.}]
Under $H_{\mathrm{CL}}(A,A^{\mathrm i},\tau_{\mathrm i})$, Priddy's
quadratic coalgebra $A^i$ maps
quasi-isomorphically to $B(A)$.  The weight spectral sequence and its
contracting homotopy prove this comparison.
\item[\textsc{Curved completed lane.}]
For a weakly curved $A_\infty$-algebra over a complete local base,
Positselski's equivalence is formed after localization at
semi-isomorphisms \cite[Theorem~6.7.1]{Positselski2018}.  A chiral
factorization realization requires its own stratumwise assembly
theorem and belongs to $H_{\mathrm{fact}}$.
\end{itemize}
The maps $A^i\to B(A)$ and
$B_X(A)\to\mathbb D_{\Ran}\mathbb D_{\Ran}B_X(A)$ organize the
intersection of these lanes.  Each archetype enters after its
weight filtration, completion, and dualizability witnesses have been
constructed.

\paragraph{Four lanes by categorical level.}
\label{par:A-four-lanes}%
\index{Theorem A!four lanes by categorical level}%
The categorical ledger has four levels:
\begin{enumerate}[label=\textsc{L\arabic*.}]
\item \label{L1-chain} \emph{Chain level;
$\ClaimStatusProvedElsewhere$.}
For augmented dg algebras and conilpotent dg coalgebras, the unit and
counit are quasi-isomorphisms
\cite[Corollary~2.3.4]{LV12}.  The $P_\infty$ form is
\cite[Proposition~11.4.4]{LV12}.
\item \label{L2-quillen} \emph{Ran $(\infty,1)$ level;
$\ClaimStatusProvedElsewhere$.}
Proposition~4.1.2 of \cite{Francis2012}, applied in the
pro-nilpotent category $(D(\Ran X),\otimes^{\ch})$, gives the enhanced
associative equivalence.
\item \label{L3-infty1} \emph{Associative factorization level;
$\ClaimStatusConjectured$.}
The package $H_{\mathrm{fact}}$ is the closure theorem carrying the
Ran equivalence to the factorization subcategories and their collision
descent data.
\item \label{L4-infty2} \emph{Properadic comparison level;
$\ClaimStatusConjectured$.}
Frontier~$A2$ asks for an enriched factorization-properad model and a
comparison functor from its operadic bar--cobar envelope.  The
presentation clause $H_{\mathrm{HR}}$
(Theorem~\ref{thm:hackney-robertson-model}, Conjectured) posits the
model structure on simplicial properads which forms one input to that
construction; the graphical-set homotopy theory of
Hackney--Robertson--Yau \cite{HRY17} is its evidence.
\end{enumerate}
Admission of an archetype to levels~\ref{L3-infty1} and
\ref{L4-infty2} is witnessed by explicit closure, completion, and
descent morphisms.  The finite-weight series of
Lemma~\ref{lem:archetype-H123-witness} supplies the locally finite
bar-weight input in $H_{\mathrm{conv}}$; the separate quadratic
acyclicity belongs to $H_{\mathrm{CL}}(A,A^{\mathrm i},\tau_{\mathrm i})$.

\paragraph{Notion of $\Eone$-chiral algebra in use.}
Five constructions occur in the surrounding literature:
(A) algebras for the reduced associative operad in
$(D(\Ran X),\otimes^{\ch})$; (B) $A_\infty$-algebras in a chiral
endomorphism operad; (C) Etingof--Kazhdan quantum vertex algebras;
(D) associative factorization algebras; and (E) ordered-Ran
factorization data.  Theorem~A begins with construction~(A).
Construction~(D) enters through $H_{\mathrm{fact}}$.

Each passage among (A)--(E) carries its own comparison morphism.  An
$A_\infty$ enhancement is governed by the internal form of
\cite[Proposition~11.4.4]{LV12}.  Quantized braided descent carries a
Drinfeld associator and a unitary quantum $R$-datum.  Ordered-to-Ran
descent carries the collision-extension maps of
$H_{\mathrm{fact}}$.  The bar object is a coalgebra over
$\operatorname{Bar}(\mathsf{Ass})$.  The Swiss--cheese action belongs
to the subsequent center pair
$(C^\bullet_{\ch}(A,A),A)$ through its own center-level theorem.

\begin{convention}[Ordered bar signs and desuspension; \ClaimStatusDefinitional]
\label{conv:A-infinity-2-ordered-bar-signs}
\index{ordered bar complex!desuspension signs}
\index{bar differential!ordered sign convention}
The cohomological convention is $|d|=+1$ and
$|s^{-1}a|=|a|-1$. For homogeneous $a_i\in\bar\cA$ the ordered bar
coalgebra is
\[
B_X^{\ord}(\cA)=T^c(s^{-1}\bar\cA)
\]
with deconcatenation coproduct. Its differential is the coderivation
whose unary and binary projections are
\[
b_1(s^{-1}a)=-s^{-1}(d_\cA a),\qquad
b_2^{(r)}(s^{-1}a\otimes s^{-1}b)
=(-1)^{|a|}s^{-1}(a_{(r)}b),
\]
and whose action on the $i$-th adjacent pair carries the prefix sign
\[
(-1)^{\sum_{q<i}(|a_q|-1)+|a_i|}.
\]
For an $A_\infty$-chiral algebra, the higher maps are extended by the
same desuspended coderivation rule. With this convention
$d_{B^{\ord}}^2=0$ is precisely the desuspended Stasheff identity
together with the Arnold relation on logarithmic collision forms.
The symmetric bar inherits this sign convention through the
$R$-twisted descent morphism of
Lemma~\ref{lem:R-twisted-descent-unitary}.

For a coaugmented dg coalgebra $C$ with reduced coproduct
$\bar\Delta c=\sum c_{(1)}\otimes c_{(2)}$, the compatible cobar
construction is
\[
 \Omega(C)=T(s\bar C),\qquad |sc|=|c|+1.
\]
On generators its differential is
\[
 d_{\Omega}(sc)
 =-s(d_Cc)-\sum(-1)^{|c_{(1)}|}
   (sc_{(1)})(sc_{(2)}).
\]
The internal term and the coproduct term both have degree~$+1$.
Equivalently, the degree-one universal twisting map
$\tau_C(c)=sc$ satisfies
$d\tau_C+\tau_C\star\tau_C=0$.
\end{convention}

%================================================================
\section{Bar-cobar reconstruction and Verdier duality: unified Theorem~A}\label{sec:ainf2-koszul-reflection}
%================================================================

\subsection{Ran reconstruction, factorization closure, and the
quadratic--Verdier refinement}

\begin{computation}[The first square-zero algebra;
\ClaimStatusProvedHere]
\label{comp:A-square-zero-bar-cobar}
Let $A=k\oplus k\epsilon$ with augmentation
$\epsilon\mapsto0$ and $\epsilon^2=0$.  Then
\[
 B(A)=T^c(s^{-1}k\epsilon).
\]
Write $e_n=[\epsilon|\cdots|\epsilon]$ for the length-$n$ bar word and
$x_n:=s e_n$ for its cobar generator.  Thus $|e_n|=-n$ and
$|x_n|=1-n$.  The bar differential vanishes,
and the cobar differential is
\[
 d x_n=\sum_{i=1}^{n-1}(-1)^{\alpha_i}x_i x_{n-i},
\]
with the sign $\alpha_i$ fixed by
Convention~\ref{conv:A-infinity-2-ordered-bar-signs}.  At fixed total
length, every term on the right has degree $2-n=|x_n|+1$.  At fixed total
length $n$, a monomial
$x_{n_1}\cdots x_{n_r}$ is indexed by the composition
$n=n_1+\cdots+n_r$, hence by a subset of the $n-1$ possible internal
cuts.  The differential adds one cut.  This identifies the
total-length-$n$ summand, for $n\geq2$, with the augmented cochain
complex of the $(n-2)$-simplex, equivalently with the Boolean cut
complex.  It is acyclic.

The total-length-$0$ and total-length-$1$ summands are generated by
$1$ and $x_1$.  The counit sends $x_1$ to $\epsilon$ and $x_n$ to zero
for $n\geq2$; moreover $x_1^2=d x_2$ up to the fixed sign.  Hence
\[
 H^\bullet(\Omega B(A))\cong k\oplus k\epsilon=A.
\]
This calculation is the first nontrivial instance of the universal
bar--cobar resolution.
\end{computation}

\begin{theorem}[Theorem~A: associative bar--cobar reconstruction in
the chiral Ran ambient; \ClaimStatusConditional]
\label{thm:koszul-reflection}
\index{Theorem A!bar-cobar reconstruction and Verdier duality}
\index{Koszul duality!Verdier-dualized chiral form}
\index{bar-cobar!reconstruction}
\textup{[Type signature: Open quadrant; chiral Ran presentation;
Beilinson level $1\leftrightarrow2$; reduced associative operad;
categorical level $(\infty,1)$; pro-nilpotent Francis--Gaitsgory
ambient $\mathcal C_X$; factorization package
$H_{\mathrm{fact}}$; named hypothesis packages consisting of the
completed chain package $H_{\mathrm{conv}}$ and the Verdier package
$H_{\mathrm{VD}}$.  Clause~\ref{KR-iv} carries the
Theorem~B package
$H_{\mathrm{CL}}(A,A^{\mathrm i},\tau_{\mathrm i})$.]}
Let $X$ be a separated finite-type scheme over a field $k$ of
characteristic zero and set
\[
 \mathcal C_X:=\bigl(D(\Ran X),\otimes^{\ch}\bigr).
\]
The following statements occupy successive domains.
\begin{enumerate}[label=\textup{(\roman*)}]
\item \label{KR-i} \emph{Enhanced Ran equivalence;
$\ClaimStatusProvedElsewhere$.}
For the reduced associative operad, enhanced bar and cobar give
mutually inverse equivalences
\[
 \operatorname{Bar}^{\mathrm{enh}}_{\mathsf{Ass}}\colon
 \operatorname{Alg}_{\mathsf{Ass}}(\mathcal C_X)
 \xrightarrow{\ \sim\ }
 \operatorname{CoAlg}^{\mathrm{dp}}_{
      \operatorname{Bar}(\mathsf{Ass})}(\mathcal C_X)
 \mathbin{:}
 \operatorname{Cobar}^{\mathrm{enh}}_{
      \operatorname{Bar}(\mathsf{Ass})}.
\]
\item \label{KR-ii} \emph{Universal reconstruction;
$\ClaimStatusProvedElsewhere$.}
For augmented unital algebras, apply clause~\ref{KR-i} to the
nonunital augmentation ideal and then restore the unit.  The unit and
counit are equivalences; on chain models they specialize to
\[
 C\xrightarrow{\ \sim\ }B\Omega C,
 \qquad
 \Omega B(A)\xrightarrow{\ \sim\ }A.
\]
Under $H_{\mathrm{conv}}$, these maps pass to the selected completed
chain realizations through the exhaustive complete filtrations and
cohomologically Mittag--Leffler inverse systems.
\item \label{KR-iii} \emph{Factorization closure;
$\ClaimStatusConjectured$.}
Under $H_{\mathrm{fact}}$, the equivalence in clause~\ref{KR-i}
restricts to the associative factorization subcategories of this
volume.  The resulting functors are denoted
\[
 \Omegach_X\colon\CoAlg_X^{\fact,\mathrm{dp}}
 \rightleftarrows
 \Alg_X^{\fact,\aug}\colon\Bbarch_X .
\]
\item \label{KR-iv} \emph{Quadratic compression;
$\ClaimStatusConditional$.}
Under the Theorem~B package
$H_{\mathrm{CL}}(A,A^{\mathrm i},\tau_{\mathrm i})$, the canonical morphism
\[
 A^i\longrightarrow B_X(A)
\]
is a quasi-isomorphism.  Equivalently,
$\Omega_X(A^i)\to A$ is a quasi-isomorphism.  Local finite
dimensionality identifies the continuous dual of $A^i$ with the
quadratic Koszul-dual algebra $A^!$.
\item \label{KR-v} \emph{Verdier refinement;
$\ClaimStatusConditional$.}
Under $H_{\mathrm{VD}}$, Verdier duality transports the bar coalgebra
to the algebra
\[
 K_X(A):=\mathbb D_{\Ran}B_X(A).
\]
The biduality morphism identifies
$B_X(A)\xrightarrow{\sim}\mathbb D_{\Ran}K_X(A)$ on this domain.
Together with clause~\ref{KR-iv}, this identifies $K_X(A)$ with the
continuous Verdier realization of $A^!$.
\end{enumerate}
\end{theorem}

\begin{proof}
Francis--Gaitsgory's Proposition~4.1.2 gives enhanced operadic
bar--cobar equivalence in every pro-nilpotent stable
symmetric-monoidal $\infty$-category.  Their proof of
Theorem~5.1.1 establishes pro-nilpotence of
$(D(\Ran X),\otimes^{\ch})$.  Applying the proposition to the reduced
associative operad proves clause~\ref{KR-i}.

Augmented unital algebras correspond to nonunital augmentation ideals
together with re-unitalization.  The equivalence of
clause~\ref{KR-i} therefore carries its unit and counit to the maps in
clause~\ref{KR-ii}.  On ordinary dg chain models this is the
bar--cobar adjunction of \cite[Theorem~2.2.9]{LV12} and the unit--counit
quasi-isomorphism of \cite[Corollary~2.3.4]{LV12}.

The closure morphisms collected in $H_{\mathrm{fact}}$ make enhanced
bar and cobar preserve the factorization equivalences on every
disjoint locus and their extensions across collision strata.  Hence
the ambient equivalence restricts as stated in clause~\ref{KR-iii}.

Under $H_{\mathrm{CL}}(A,A^{\mathrm i},\tau_{\mathrm i})$, the universal
twisting morphism $A^i\to A$ is Koszul.  The twisting-morphism fundamental theorem
\cite[Theorem~2.3.2]{LV12} identifies this condition with either
$A^i\to B(A)$ or $\Omega(A^i)\to A$ being a quasi-isomorphism.  This
proves clause~\ref{KR-iv}.

Finally, $H_{\mathrm{VD}}$ supplies a dualizable bar object and the
monoidal Verdier comparison maps.  Verdier duality transports its
coalgebra structure to $K_X(A)$, and Verdier biduality gives the
displayed equivalence in clause~\ref{KR-v}.
\end{proof}

\subsection{The single diagram and the single equation}

The ambient equivalence is captured by one diagram:
\[
\begin{tikzcd}[column sep=large,row sep=large]
\operatorname{CoAlg}^{\mathrm{dp}}_{\operatorname{Bar}(\mathsf{Ass})}
 (\mathcal C_X)
 \arrow[r,bend left=15,"\operatorname{Cobar}^{\mathrm{enh}}"]
 \arrow[r,phantom,"\rotatebox{90}{$\dashv$}"description]
&
\operatorname{Alg}_{\mathsf{Ass}}(\mathcal C_X)
 \arrow[l,bend left=15,"\operatorname{Bar}^{\mathrm{enh}}"]
\end{tikzcd}
\qquad
\mathcal C_X=(D(\Ran X),\otimes^{\ch}).
\]
Its two equations are
\[
\boxed{
\operatorname{Cobar}^{\mathrm{enh}}
\operatorname{Bar}^{\mathrm{enh}}\simeq\mathrm{id}_{\Alg},
\qquad
\operatorname{Bar}^{\mathrm{enh}}
\operatorname{Cobar}^{\mathrm{enh}}\simeq\mathrm{id}_{\CoAlg^{\mathrm{dp}}}.
}
\]
Under $H_{\mathrm{fact}}$ the same equations restrict to
$\Omegach_X$ and $\Bbarch_X$ on the associative factorization
subcategories.

\begin{remark}[Reconstruction and duality are separate operations]
\label{rem:reconstruction-duality-separation}
\index{bar-cobar!reconstruction versus Verdier duality}
Two functors leave the bar object in different directions.  The
counit
$\Omegach_X\Bbarch_X(A)\to A$ gives reconstruction.  Verdier duality
gives the Koszul-dual lane
$K_X(A)=\mathbb D_{\Ran}\Bbarch_X(A)$.  Their codomains and hypothesis
packages determine their respective type signatures.
\end{remark}

\begin{corollary}[Typed comparison of reconstruction and Verdier
duality;
\ClaimStatusProvedHere]
\label{cor:no-dual-from-barcobar-counit}
\index{bar-cobar!comparison with Verdier duality}
\index{Koszul duality!requires Verdier branch}
\index{self-duality theorem!required for identifying reconstruction with duality}
Type signature: \textup{(}Open quadrant, reconstruction/Verdier
comparison lane, Beilinson level \(1\leftrightarrow2\),
pro-nilpotent chiral Ran ambient; packages
$H_{\mathrm{fact}}\cup H_{\mathrm{VD}}$\textup{)}.
Let $A$ lie in the factorization and Verdier domains of
Theorem~\ref{thm:koszul-reflection}.  Reconstruction and Verdier
duality give
\[
 \varepsilon_A\colon R_X(A)=\Omegach_X\Bbarch_X(A)\xrightarrow{\sim}A,
 \qquad
 K_X(A)=\mathbb D_{\Ran}\Bbarch_X(A).
\]
Composition with $\varepsilon_A$ induces an equivalence of mapping
spaces
\[
 \operatorname{Map}(A,K_X(A))
 \xrightarrow{\ \sim\ }
 \operatorname{Map}(R_X(A),K_X(A)).
\]
Hence an equivalence $R_X(A)\simeq K_X(A)$ is precisely the composite
of reconstruction with a self-duality equivalence
\[
\sigma_A\colon A\xrightarrow{\sim}K_X(A).
\]
Under $H_{\mathrm{CL}}(A,A^{\mathrm i},\tau_{\mathrm i})$ and local finite
duality, the target
$K_X(A)$ identifies with the continuous realization of $A^!$.  A
strict model is obtained from a specified formality equivalence
$K_X(A)\simeq A^!$.
\end{corollary}

\begin{proof}
The morphism $\varepsilon_A$ is an equivalence by
Theorem~\ref{thm:koszul-reflection}\ref{KR-ii}.  Precomposition with
an equivalence induces an equivalence of mapping spaces.  Under this
equivalence, $\sigma_A$ corresponds to
$\sigma_A\circ\varepsilon_A$.  Clauses~\ref{KR-iv} and~\ref{KR-v}
supply the quadratic and Verdier identifications of the target.
\end{proof}

\begin{remark}[Four named morphisms in Theorem~A]
\label{rem:four-named-morphisms}
\index{Koszul reflection!four named morphisms}
Four canonical morphisms organize the theorem:
\begin{itemize}
\item $\eta_C\colon C\xrightarrow{\sim}B\Omega(C)$ is the
coalgebra-side unit in the Ran ambient;
\item $\varepsilon_A\colon\Omega B(A)\xrightarrow{\sim}A$ is the
algebra-side counit in the Ran ambient;
\item $q_A\colon A^i\to B(A)$ is the quadratic comparison under
$H_{\mathrm{CL}}(A,A^{\mathrm i},\tau_{\mathrm i})$;
\item $\delta_{B(A)}\colon
B(A)\xrightarrow{\sim}\mathbb D_{\Ran}\mathbb D_{\Ran}B(A)$ is
Verdier biduality under $H_{\mathrm{VD}}$.
\end{itemize}
The factorization closure package transports these morphisms to the
factorization subcategories.
\end{remark}

\begin{proposition}[Module transport along the bar--cobar counit;
\ClaimStatusProvedHere]
\label{prop:A-module-completion-firewall}
\index{Theorem A!module and completion ambient}
\index{coderived category!bar-cobar continuation}
Let $A\in\operatorname{Alg}_{\mathsf{Ass}}(\mathcal C_X)$ and put
$R_X(A)=\operatorname{Cobar}^{\mathrm{enh}}
\operatorname{Bar}^{\mathrm{enh}}(A)$.  Extension and restriction
along the counit give mutually inverse equivalences
\[
\operatorname{Mod}_{R_X(A)}(\mathcal C_X)
\simeq
\operatorname{Mod}_{A}(\mathcal C_X).
\]
Under $H_{\mathrm{fact}}$ this equivalence restricts to the
factorization-module categories.  Under a self-duality equivalence
$\sigma_A\colon A\xrightarrow{\sim}K_X(A)$ it further gives
\[
\operatorname{Mod}_{R_X(A)}^{\fact}
\simeq
\operatorname{Mod}_{A}^{\fact}
\simeq
\operatorname{Mod}_{K_X(A)}^{\fact}.
\]
\end{proposition}

\begin{proof}
The counit is an equivalence by
Theorem~\ref{thm:koszul-reflection}\ref{KR-ii}.  The functoriality of
module objects carries an equivalence of associative algebras to an
equivalence of their module categories.  The packages
$H_{\mathrm{fact}}$ and $\sigma_A$ supply the two displayed
restrictions and transports.
\end{proof}

\begin{remark}[Operadic levels of the bar object and the center pair]
\label{rem:bar-over-chirAss-shriek}

\index{SC-chtop!derived-centre action}
The bar object $\Bbarch_X(A)$ is a coalgebra over
$\operatorname{Bar}(\mathsf{Ass})$, with deconcatenation coproduct
and degree-$1$ bar differential.  The derived center pair
$(C^\bullet_{\ch}(A,A),A)$ carries its brace structure at the next
operadic level.  A Swiss--cheese action on this pair is supplied by
the corresponding open--closed center theorem.  Thus the bar object
serves as the level-$2$ resolution and the center pair as the
level-$3$ open--closed object.
\end{remark}

\subsection{Locally finite positive weights and bar finiteness}
\label{subsec:archetype-witness-H123}

Positive weight bounds the tensor length contributing to a fixed bar
weight.  This elementary observation is the complete finiteness input
needed here.

For the square-zero algebra of
Computation~\ref{comp:A-square-zero-bar-cobar}, assign weight~$1$ to
$\epsilon$.  Its length-$p$ bar word has total weight~$p$, so the
weight-$w$ reduced bar is the one-dimensional span of $e_w$.  The
general statement replaces this single word by a finite sum over
compositions of~$w$.

\begin{lemma}[Positive local finiteness passes to the reduced bar;
\ClaimStatusProvedHere]
\label{lem:archetype-H123-witness}
\index{bar construction!positive-weight local finiteness}%
\textup{[Type signature: Open quadrant; weight-graded chain
presentation; Beilinson level~$2$; hypothesis package: augmentation,
positive weight, and finite-dimensional weight fibers.]}
Let $A=k\mathbf 1\oplus\bar A$ be an augmented dg algebra and suppose
\[
 \bar A=\bigoplus_{w\geq w_0}\bar A_w,
 \qquad w_0\geq1,
 \qquad \dim_k\bar A_w<\infty.
\]
Assume that $d_A$ has weight zero and that multiplication has additive
weight.  Then the reduced bar
$\bar B(A)=\bigoplus_{p\geq1}(s^{-1}\bar A)^{\otimes p}$ is graded by
total weight, its differential preserves that grading, and every
weight space $\bar B(A)_w$ is finite-dimensional.  Writing
$d_w=\dim_k\bar A_w$ and $P_A(q)=\sum_{w\geq w_0}d_wq^w$, one has
\[
\dim_k \bar B(A)_w
 \;=\;
 \sum_{p\geq 1}\;
 \sum_{w_1+\cdots+w_p=w,\,w_i\geq w_0}\;
 d_{w_1}\cdots d_{w_p}
 \;=\; [q^w]\,\frac{P_A(q)}{1-P_A(q)}.
\]
The same conclusion holds for a weight-graded factorization algebra
after application of any exact weight-preserving symmetric-monoidal
fiber functor to $\mathsf{Ch}(k)$.
\end{lemma}

\begin{proof}
Every tensor in total weight $w$ has length
$p\leq\lfloor w/w_0\rfloor$.  For fixed $p$, the set of compositions
$w=w_1+\cdots+w_p$ with $w_i\geq w_0$ is finite.  Each corresponding
tensor product
$\bar A_{w_1}\otimes\cdots\otimes\bar A_{w_p}$ is finite-dimensional,
so their finite direct sum is finite-dimensional.  The internal
differential preserves every $w_i$, while the multiplication term
replaces adjacent weights $w_i,w_{i+1}$ by $w_i+w_{i+1}$; hence the bar
differential preserves total weight.  Summing over tensor length gives
$\sum_{p\geq1}P_A(q)^p=P_A(q)/(1-P_A(q))$.  An exact fiber functor
preserves the displayed direct sums, tensor products, and differential.
\end{proof}

\begin{remark}[Weight and length filtrations]
\label{rem:H3-not-finite-bar-depth}
\index{bar construction!weight and length filtrations}%
The length-$p$ summand is $(s^{-1}\bar A)^{\otimes p}$ and may carry
infinitely many total weights.  The weight-$w$ summand receives
contributions from the finite range
$1\leq p\leq\lfloor w/w_0\rfloor$.  Thus the direct-sum length
filtration supplies $H_{\mathrm{len}}$, while the positive grading and
finite-dimensional fibers supply finite bar-weight pieces.
Augmentation completeness and Koszul acyclicity remain the separate
maps in $H_{\mathrm{aug}}$ and
$H_{\mathrm{CL}}(A,A^{\mathrm i},\tau_{\mathrm i})$.
\end{remark}

\begin{remark}[Class-$\mathsf{M}$ completion problem;
\ClaimStatusConditional]
\label{rem:class-M-pro-conilpotent}
\index{class M!pro-conilpotent completion}%
\index{pro-conilpotent completion!Virasoro and W-algebras}%
Let $A$ be a positive-energy Virasoro or principal $\mathcal W_N$
factorization model with finite-dimensional weight fibers.  The
weight-completed object is
\[
\widehat A^{\mathrm{wt}}=\varprojlim_N A/A_{\geq N}
 \;\in\; \Alg^{\fact,\aug,\comp}_X,
\]
while augmentation completion is
\[
\widehat A^{I}=\varprojlim_q A/\bar A^{\,q}.
\]
The class-$\mathsf M$ application of Theorem~A requires an explicit
cofinal comparison
$\gamma_A\colon\widehat A^{I}\xrightarrow{\sim}\widehat A^{\mathrm{wt}}$,
the equivalence $A\xrightarrow{\sim}\widehat A^{I}$ of
$H_{\mathrm{aug}}$, the pro-coradical Mittag--Leffler maps of
$H_{\mathrm{ML}}$, and the acyclicity of
$A\otimes_{\tau_{\mathrm i}}A^i$ from
$H_{\mathrm{CL}}(A,A^{\mathrm i},\tau_{\mathrm i})$.
\end{remark}

\begin{remark}[Five conditional applications; \ClaimStatusConditional]
\label{rem:five-archetype-K-squared}
\index{bar-cobar reconstruction!conditional family inputs}%
Lemma~\ref{lem:archetype-H123-witness} applies to a family after its
augmentation, positive grading, and finite-dimensional weight fibers
have been constructed.  The full Ran, factorization, completion,
quadratic, and duality application lies in the pro-nilpotent chiral
Ran ambient and uses $H_{\mathrm{fact}}$, $H_{\mathrm{aug}}$, $H_{\mathrm{ML}}$,
$H_{\mathrm{CL}}(A,A^{\mathrm i},\tau_{\mathrm i})$, and $H_{\mathrm{VD}}$.
For the five named families the remaining inputs are these:
\begin{itemize}
\item[$\mathsf G$.]  For a Heisenberg or free-field presentation,
construct the vacuum augmentation as a factorization morphism and the
quadratic twisting complex
$A\otimes_{\tau_{\mathrm i}}A^i$.  Its positive-weight acyclicity and the
completion maps give the $\mathsf G$ application.
\item[$\mathsf L$.]  For $V_k(\fg)$, construct the factorization bar at
the chosen level, the cofinal map from augmentation completion to the
chosen level/weight completion, and a contraction of
$A\otimes_{\tau_{\mathrm i}}A^i$.  At $k=-h^\vee$ the corresponding statement
uses a separately specified completed resolution together with its
maps $\iota_A,p_A,h_A$.
\item[$\mathsf C$.]  For $\beta\gamma_\lambda$, choose a charge sector
whose conformal weights are bounded below and finite-dimensional,
then construct the factorization collision maps and the quadratic
twisting contraction on that sector.
\item[$\mathsf M$.]  For Virasoro and principal $\mathcal W_N$, supply
the comparison $\gamma_A$ of
Remark~\ref{rem:class-M-pro-conilpotent}, the pro-coradical
Mittag--Leffler tower, and a chain map from the relevant Felder or BRST
resolution to $A\otimes_{\tau_{\mathrm i}}A^i$ carrying its contraction.
\item[$\mathsf B$.]  For a Mukai/Hall presentation, choose a
positive-energy cone or charge truncation with finite-dimensional
weight fibers, construct its factorization multiplication and
augmentation, and give the quadratic twisting contraction in that
completed sector.
\end{itemize}
Each item is a concrete input list for
Theorem~\ref{thm:koszul-reflection}\ref{KR-iii}--\ref{KR-v}.
\end{remark}

\subsection{Rescaling a deformation retract}
\label{subsec:A-universal-chain-homotopy}

A scalar rescales a deformation retract by an elementary identity.
Its identification with a trace or conductor requires a family-level
chain comparison.

\begin{convention}[Scalar invariant and chain-level rescaling]
\label{conv:A-two-kappa-shriek}
\index{Verdier duality!at level 1 vs level 2}%
The algebra-level scalar invariant is
\[
 K^\kappa(\cA):=\kappa(\cA)+\kappa(\cA^!).
\]
It is defined from the chosen scalar theory on a Verdier--Koszul pair.
A chain-level rescaling is a separate datum: for a deformation retract
$(\iota_A,p_A,h_A)$ of a complex $(C_A,d_A)$, it is a scalar
$\nu_A\in k$ occurring in the identity
\[
 d_Ah_A+h_Ad_A
 =\nu_A(\mathrm{id}_{C_A}-\iota_Ap_A).
\]
A chain map relating the contraction to the scalar trace promotes this
datum to the theorem $\nu_A=K^\kappa(A)$.
\end{convention}

\begin{proposition}[Rescaling identity for a deformation retract;
\ClaimStatusProvedHere]
\label{prop:A-universal-chain-homotopy}
\index{chain-homotopy!scalar rescaling}%
\textup{[Type signature: chain presentation; deformation-retract
level; hypothesis package: a scaled homotopy identity with invertible
scalar.]}
Let $(C,d)$ and $(H,d_H)$ be complexes, let
$\iota\colon H\to C$ and $p\colon C\to H$ be chain maps with
$p\iota=\mathrm{id}_H$, and let $h\colon C\to C[-1]$.  Suppose that
$\nu\in k^\times$ and
\[
 dh+hd=\nu(\mathrm{id}_C-\iota p).
\]
Then $\widetilde h=\nu^{-1}h$ satisfies
\[
 d\widetilde h+\widetilde h d=\mathrm{id}_C-\iota p.
\]
The side conditions $h\iota=0$, $ph=0$, and $h^2=0$ pass to
$\widetilde h$.  Conversely, every contracting homotopy $h_0$ gives a
scaled identity with $h=\nu h_0$.
\end{proposition}

\begin{proof}
Scalar linearity gives
\[
d(\nu^{-1}h)+(\nu^{-1}h)d
=\nu^{-1}(dh+hd)
=\mathrm{id}_C-\iota p.
\]
Multiplication by $\nu^{-1}$ preserves the three side conditions.
Starting from $dh_0+h_0d=\mathrm{id}_C-\iota p$ and setting
$h=\nu h_0$ gives the converse identity.
\end{proof}

\begin{remark}[Family comparison problems for scalar rescaling;
\ClaimStatusConditional]
\label{rem:A-h-three-paths}
\index{chain-homotopy!family comparison problems}%
For a family $F$, put
$C_F=\Omegach_X\Bbarch_X(A_F)$ and $p_F=\varepsilon_{A_F}$.  A proposed
scalar $\nu_F$ becomes a chain-level rescaling after constructing a
chain map $\iota_F\colon A_F\to C_F$ and an operator
$h_F\colon C_F\to C_F[-1]$ satisfying
\[
p_F\iota_F=\mathrm{id}_{A_F},
\qquad
d_Fh_F+h_Fd_F
=\nu_F(\mathrm{id}_{C_F}-\iota_Fp_F).
\]
The five scalar expressions lead to five separate construction
problems.
\begin{itemize}
\item[$\mathsf G$: $\nu_{\mathsf G}=1$.]
Construct $\iota_{\mathsf G}$ and $h_{\mathsf G}$ on the actual
factorization bar--cobar complex from the Fock/Koszul resolution.
\item[$\mathsf L$: $\nu_{\mathsf L}=2(k+h^\vee)$.]
Localize the level-parameter ring at $k+h^\vee$.  Construct a chain
map from the chosen Garland--Lepowsky or Sugawara resolution to
$C_{\mathsf L}$, together with $\iota_{\mathsf L}$ and
$h_{\mathsf L}$ satisfying the displayed scaled identity.
\item[$\mathsf C$: $\nu_{\mathsf C}=1$.]
On the selected finite charge sector, construct the chain map from the
free $\beta\gamma$ Koszul resolution to $C_{\mathsf C}$ and transport
its deformation retract.
\item[$\mathsf M$: $\nu_{\mathsf M}=c(5c+22)/10$.]
Localize the central-charge parameter ring at $\nu_{\mathsf M}$.
Construct a chain map from the Felder/BRST resolution to
$C_{\mathsf M}$ and operators
$\iota_{\mathsf M},h_{\mathsf M}$ satisfying the scaled identity.
\item[$\mathsf B$: $\nu_{\mathsf B}=8$.]
Construct a chain map from the completed Mukai/Hall resolution to
$C_{\mathsf B}$ and operators $\iota_{\mathsf B},h_{\mathsf B}$
satisfying the scaled identity.
\end{itemize}
An equality $\nu_F=K^\kappa(A_F)$ consists additionally of the chain
map in Convention~\ref{conv:A-two-kappa-shriek} comparing the
contraction datum with the scalar trace complex.
\end{remark}

\begin{remark}[The chain comparison required between Theorems~A, C,
and D]
\label{rem:A-h-bridge}
\index{chain rescaling!bridge to Theorems C, D}%
Let $(C_A,d_A)$ be the complex used in Theorem~A.  A scalar
$\nu_A\in k^\times$ rescales a homotopy after one has constructed
maps $\iota_A,p_A$ and an operator $h_A$ such
that
\[
 d_Ah_A+h_Ad_A
 =\nu_A\bigl(\mathrm{id}_{C_A}-\iota_Ap_A\bigr).
\]
Then $\nu_A^{-1}h_A$ gives the corresponding
contraction.  A relation with Theorems~C and~D therefore requires an
explicit chain map that identifies $\nu_A$ with the scalar
invariant occurring in the census or Hodge obstruction.  The five
family expressions
$\{1,2(k+h^\vee),1,c(5c+22)/10,8\}$ define candidate values for
$\nu_A$; the
construction of this chain map is the comparison problem linking the
three theorems.
\end{remark}

%================================================================
\section{The eight standard facets of Theorem~A}\label{sec:ainf2-eight-corollaries}
%================================================================

The eight standard facets of Theorem~A are recovered as corollaries of
the unified Koszul reflection
Theorem~\ref{thm:koszul-reflection}.

\begin{corollary}[Twisting morphism $\tau\colon\Bbarch_X(\cA)\to\cA$
is a Maurer--Cartan element; \ClaimStatusConditional]
\label{cor:eight-cor-twisting-morphism}
\index{twisting morphism!as Maurer-Cartan}
The universal twisting morphism
$\tau\colon\Bbarch_X(\cA)\to\cA$ associated to the algebra-side counit
$\psi_\cA\colon\Omegach_X\Bbarch_X(\cA)\to\cA$ of the Koszul
reflection is a Maurer--Cartan element of the convolution
$L_\infty$-algebra
$\fg^{\Eone}_\cA=\Hom^{\fact}_X(\Bbarch_X(\cA),\cA)$:
\[
d\tau+\tfrac{1}{2}[\tau,\tau]=0.
\]
\end{corollary}

\begin{proof}
By Theorem~\ref{thm:koszul-reflection}\ref{KR-i}, the counit
of the adjunction $\Omegach_X\dashv\Bbarch_X$ encodes, by adjunction,
the universal twisting
morphism $\tau\colon\Bbarch_X(\cA)\to\cA$ (LV12 Theorem~2.2.9
specialised to the chiral ambient). The Maurer--Cartan condition
$d\tau+\tfrac{1}{2}[\tau,\tau]=0$ is exactly the compatibility
between the bar differential on $\Bbarch_X(\cA)$ and the
multiplication on~$\cA$, which is the chain-level statement
that~$\psi_\cA$ is a chain map.
\end{proof}

\begin{corollary}[Universal reconstruction in the enhanced Ran
ambient; \ClaimStatusProvedElsewhere]
\label{cor:eight-cor-counit-qi}
\index{counit!universal reconstruction in the Ran ambient}
\textup{[Type signature: Open quadrant; enhanced associative Ran
presentation; Beilinson levels (1\leftrightarrow2); pro-nilpotent
Francis--Gaitsgory ambient; reduced associative operad.]}
For every augmented algebra
$\cA\in\operatorname{Alg}^{\mathrm{aug}}_{mathsf{Ass}}
(D(\operatorname{Ran}X),\otimes^{\mathrm{ch}})$, the enhanced counit
\[
 \varepsilon_\cA\colon
 \Omega_X^{\mathrm{enh}}B_X^{\mathrm{enh}}(\cA)
 \xrightarrow{\ \sim\ }\cA
\]
is an equivalence.
\end{corollary}

\begin{proof}
This is Theorem~\ref{thm:koszul-reflection}\ref{KR-ii}, obtained by
applying Francis--Gaitsgory's pro-nilpotent enhanced bar--cobar
equivalence to the augmentation ideal and restoring the unit.
\end{proof}

\begin{remark}[Factorization and chain realizations]
\label{rem:eight-cor-counit-factorization-realization}
Under (H_{\mathrm{fact}}), the displayed equivalence restricts to the
chosen factorization subcategories.  Under (H_{\mathrm{conv}}), its
completed chain realization is a quasi-isomorphism.  Quadratic
Koszulness is the further comparison
(q_\cA\colon\cA^{\mathrm i}\to B_X(\cA)) of
Theorem~\ref{thm:koszul-reflection}\ref{KR-iv}.
\end{remark}

\begin{corollary}[Unit weak equivalence on $\Conil(X)$;
\ClaimStatusConditional]
\label{cor:eight-cor-unit-weq}
\index{unit!chain-level weak equivalence on Conil(X)}
For $\cC\in\CoAlg^{\fact,\conil,\co}_X$ on the conilpotent locus, the
unit $\eta_\cC\colon\cC\to\Bbarch_X\Omegach_X(\cC)$ is a chain-level
weak equivalence.
\end{corollary}

\begin{proof}
Filter $\cC$ by the coradical filtration. The associated graded of
$\Bbarch_X\Omegach_X(\cC)$ is the classical bar construction on the
free algebra generated by the graded pieces of $\cC$, and the map
$\gr\,\eta_\cC$ is the unit of the acyclic twisting morphism for
$(\Ass,\Ass^!)$ \cite[Theorem~2.4.1]{LV12}. Conilpotence makes the
filtration exhaustive, completeness makes it separated, and
$H_{\mathrm{ML}}$ gives Mittag--Leffler convergence weight-by-weight. Hence the
filtered cone of $\eta_\cC$ is acyclic, so the unit is a chain-level
weak equivalence on the conilpotent-complete locus.
\end{proof}

\begin{corollary}[Twisted tensor acyclicity;
\ClaimStatusConditional]
\label{cor:eight-cor-twisted-tensor}
\index{twisted tensor product!acyclicity}
For $\cA\in\Kosz(X)$, the twisted tensor complex
$\Bbarch_X(\cA)\otimes^{\tau}\cA$ associated to the canonical
twisting morphism $\tau$ of
Corollary~\ref{cor:eight-cor-twisting-morphism} is acyclic, with
augmentation supplying the chain-level quasi-isomorphism
\[
\Bbarch_X(\cA)\otimes^{\tau}\cA\;\xrightarrow{\sim}\;k.
\]
\end{corollary}

\begin{proof}
The twisted tensor complex computes $\Tor^{\Omegach_X\Bbarch_X(\cA)}_{*}(k,k)$
via the bar resolution; under
Corollary~\ref{cor:eight-cor-counit-qi},
$\Omegach_X\Bbarch_X(\cA)\simeq\cA$, so the Tor reduces to
$\Tor^{\cA}_{*}(k,k)=H_{*}(\Bbarch_X(\cA))$, which by definition of
$\Kosz(X)$ is concentrated in bar degree~$0$ at the augmentation. The
augmentation map then realises the twisted tensor complex as the
augmentation chain complex, which is acyclic
(\cite[Theorem~2.4]{LV12} specialised to the chiral ambient).
\end{proof}

\begin{corollary}[Bar concentrated in weight~$1$ at $g=0$ on
class~$\mathsf{G}$; \ClaimStatusConditional]
\label{cor:eight-cor-bar-weight-1}
\index{bar complex!concentration in weight 1}
\index{class G!bar concentration}
For $\cA$ in class~$\mathsf{G}$ (PBW $E_2$-collapse class), lying in
the genus-zero finite-type PBW window of
Definition~\ref{def:ftm-g0-pbw-window}, $\Bbarch_X(\cA)$ has
cohomology concentrated in bar weight~$1$:
\[
H^{*}(\Bbarch_X(\cA))=H^{1}(\Bbarch_X(\cA))\;[\text{shifted}\,]\,\quad
H^{i}(\Bbarch_X(\cA))=0\;\text{for}\;i\neq 1.
\]
At $g\geq 1$, weight-$1$ concentration occupies the
class~$\mathsf{G}$ locus.  In classes~$\mathsf{L}$, $\mathsf{C}$,
$\mathsf{M}$, the fiberwise curvature
$d^{2}_{\mathrm{fib}}=\kappa\cdot\omega_g$ produces the higher-weight
terms.
\end{corollary}

\begin{proof}
By Theorem~\ref{thm:koszul-reflection}\ref{KR-ii}, reconstruction on
$\Kosz(X)$ is equivalent (via the PBW spectral sequence of the bar
filtration) to the $E_2$-collapse of the
bar--cobar spectral sequence at $E_2$-page. On the genus-zero
finite-type PBW window of Definition~\ref{def:ftm-g0-pbw-window} ---
and only there --- the $E_2$-collapse is
exactly the hub-and-spoke TFAE of
Theorem~\ref{thm:ftm-seven-fold-tfae-via-hub-spoke}, which gives bar
concentration in weight~$1$ at $g=0$ for any class
$\cA\in\Kosz(X)$ lying in that window; $E_2$-collapse outside the
window is not delivered by that theorem and is an open problem.
At $g\geq 1$, the fiberwise differential
$d_{\mathrm{fib}}^{2}=\kappa\cdot\omega_g$ identifies weight-$1$
concentration with the zero-curvature locus $\kappa=0$, namely
class~$\mathsf{G}$.
\end{proof}

\begin{corollary}[SC-formality criterion on class~$\mathsf{G}$;
\ClaimStatusConditional]
\label{cor:eight-cor-sc-formality}
\index{SC-formality!class G criterion}
\index{class G!SC-formal}
For $\cA$ in class~$\mathsf{G}$ lying in the genus-zero finite-type
PBW window of Definition~\ref{def:ftm-g0-pbw-window}, the convolution
$L_\infty$-algebra
$\fg^{\Eone}_\cA=\Hom^{\fact}_X(\Bbarch_X(\cA),\cA)$ is SC-formal: all
higher Stasheff--Massey operations $m_k$ for $k\geq 3$ vanish on
cohomology, and the $A_\infty$-quasi-isomorphism inverting
$\varepsilon_\cA$ collapses to a strict isomorphism.  In
classes~$\mathsf{L}$, $\mathsf{C}$, $\mathsf{M}$, the higher
operations realize the shadow-tower invariants $S_k$ and retain the
$A_\infty$ deformation.
\end{corollary}

\begin{proof}
By Theorem~\ref{thm:koszul-reflection}\ref{KR-ii}, SC-formality is
equivalent to a strict isomorphism realizing the chain-level
equivalence $\Omegach_X\Bbarch_X(\cA)\simeq\cA$.
The shadow tower $S_k$ measures the $A_\infty$-corrections to
$\varepsilon_\cA$; vanishing of $S_k$ for $k\geq 3$ is the
PBW-collapse-to-strict criterion, which by the hub-and-spoke TFAE
(Theorem~\ref{thm:ftm-seven-fold-tfae-via-hub-spoke}, spoke~7, valid
on the genus-zero finite-type PBW window of
Definition~\ref{def:ftm-g0-pbw-window}) is the
class~$\mathsf{G}$ condition exactly on that window.
\end{proof}

\begin{corollary}[$R$-twisted $\Sigma_n$-descent (unitary $R$)
$\barB^{\Sigma}(\cA)\simeq B^{\ord}(\cA)^{R\text{-}\Sigma_n}$;
\ClaimStatusConditional]
\label{cor:eight-cor-R-descent}
\index{R-twisted descent!corollary form}
For $\cA$ an $\Eone$-chiral algebra with meromorphic quantum
$R$-matrix~$R(z)$, with checked braiding
$\check R(z)=\tau\circ R(z)$, satisfying the quantum Yang--Baxter
equation (Lemma~\ref{lem:R-twisted-descent-unitary}~\ref{R1}) and unitarity
$R^{21}(-z)R^{12}(z)=\id$
(Lemma~\ref{lem:R-twisted-descent-unitary}~\ref{R2}), the restriction of
$\barB^{\Sigma}_n(\cA)$ to $\Conf_n(X)$ is the $R$-twisted descent of
the ordered bar along the principal torsor
$\pi_n^\circ\colon\Conf_n^{\ord}(X)\to\Conf_n(X)$.  Under
$H_{\mathrm{fact}}^{R}$ this descent extends to the cardinality-$n$
Ran stratum and gives
\[
\barB^{\Sigma}_n(\cA)\;\simeq\;
\bigl(B^{\ord}_n(\cA)\bigr)^{\Sigma_n\text{-coinv},\;L_R\text{-twisted}}.
\]
\end{corollary}

\begin{proof}
Lemma~\ref{lem:R-twisted-descent-unitary} constructs the equivariant
local system and proves descent on the open torsor.  The clause
$H_{\mathrm{fact}}^{R}$ supplies the compatible nearby-cycle
extension.
\end{proof}

\begin{corollary}[Conditional properadic envelope of the Ran
equivalence; \ClaimStatusConditional]
\label{cor:eight-cor-properad-equiv}
\index{Francis-Gaitsgory equivalence!properad level}
Assume $H_{\mathrm{fact}}$, $H_{\mathrm{conv}}$, and
$H_{\mathrm{prop}}$.  The enhanced bar--cobar equivalence of
Theorem~\ref{thm:koszul-reflection}\ref{KR-i} extends by
connected-graph substitution to a derived adjoint equivalence
\[
\Bbarch_X\;:\;\FactProp^{\aug}_{\langle\chirAss\rangle}(X)
\;\rightleftarrows\;
\CoFactProp^{\conil,\comp}_{\langle\chirAss\rangle}(X)
\;:\;\Omegach_X
\]
on the conilpotent-complete properadic envelope generated by the
chiral associative bar--cobar pair.  Pullback to a marked point
recovers the ordinary associative bar--cobar equivalence
\cite[Corollary~2.3.4]{LV12}.
\end{corollary}

\begin{proof}
Detailed proof in \S\ref{sec:ainf2-main} as
Theorem~\ref{thm:A-infinity-2}.
\end{proof}

These eight corollaries, taken together, account for the theorem content
of the standard Theorem~A facets. The unified reconstruction theorem
Theorem~\ref{thm:koszul-reflection} is the source of these facets;
each facet specialises to a
single facet of the adjunction, reconstruction counit, or
Verdier-dualized Koszul-duality functor.

%================================================================
\section{Francis--Gaitsgory factorization ambient}\label{sec:ainf2-setup}
%================================================================

Francis--Gaitsgory place enhanced operadic bar and cobar in the stable
presentable symmetric-monoidal category
$(D(\Ran X),\otimes^{\ch})$.  The factorization restriction selects
the coalgebras whose disjoint-locus structure maps are equivalences.
The conditional packages below record the additional closure and
model-categorical input needed for internal properads and completed
inverse limits.  Ordered descent begins on the principal
$\Sigma_n$-cover of the open configuration locus; its collision
extension is the clause $H_{\mathrm{fact}}^{R}$.

Let $X$ be a smooth curve over a field~$k$ of characteristic~$0$.
Write $\Ran(X)$ for the Ran space of~$X$ (the colimit of finite powers
$X^n$ along surjective finite-order maps), and $\Ran^{\ord}(X)$ for the
ordered Ran space (the same colimit taken over strictly increasing
index maps on labelled finite sets).

\begin{definition}[Ran factorization ambient;
\ClaimStatusDefinitional]
\label{def:fg-ambient}
\index{Francis--Gaitsgory ambient}
\index{factorization algebra!infinity-category}
The \emph{Ran factorization ambient} has sheaf-level base
\[
\Fact_{\Ran}(X)
\;=\;
\Alg^{\fact}(\Shv^{\DR}(\Ran(X)))
\]
of factorization algebras on $\Ran(X)$ valued in de~Rham sheaves, with
monoidal operation given by the Francis--Gaitsgory chiral product
\cite{Francis2012}.  Throughout this chapter $\Fact(X)$ denotes a
chosen symmetric-monoidal model enhancement satisfying
$H_{\Fact}(X)$ of Proposition~\ref{prop:fg-ambient-properties}.
\end{definition}

\begin{remark}[Monoidal factorization realization]
\label{rem:why-fact}
\index{chiral pseudo-tensor!Ran symmetric-monoidal realization}
The Beilinson--Drinfeld chiral pseudo-tensor structure on $\Dmod(X)$
defines a sequence of multilinear operations
\[
\otimes^{\ch}_I \colon \prod_{i\in I} \Dmod(X) \to \Dmod(X),\qquad
|I|<\infty,
\]
indexed by finite sets~$I$ \cite[\S3.4]{BD04}.  Its Ran realization is
the chiral symmetric-monoidal structure of Francis--Gaitsgory.  The
unit and the disjoint-union product are encoded by the factorization
structure, and $H_{\Fact}(X)$ chooses a model presentation on which
ordinary enriched transfer arguments may be applied.
\end{remark}

\begin{proposition}[Conditional factorization ambient package;
\ClaimStatusConditional]
\label{prop:fg-ambient-properties}
Let $\Fact(X)$ be the enhancement named in
Definition~\ref{def:fg-ambient}. The hypothesis package
$H_{\Fact}(X)$ consists of the following clauses.
\begin{enumerate}[label=\textup{(\alph*)}]
\item The Ran sheaf ambient $\Shv^{\DR}(\Ran(X))$ is stable and
presentable, and $\Fact_{\Ran}(X)$ is the factorization-algebra lane on
that ambient.
\item The Francis--Gaitsgory star-product extends to a colimit-compatible
symmetric monoidal operation on the chosen enhancement.
\item The factorization unit, represented sheaf-theoretically by the
constant factorization algebra on $\cO_X$, gives a strict unit in the
enhancement; the Beilinson--Drinfeld chiral unit maps to this unit by
the sheafified factorization comparison.  The two units belong to the
pseudo-tensor and factorization presentations, respectively.
\item The model enhancement is cofibrantly generated and supports the
enriched collections from which internal operads and properads are
formed.
\end{enumerate}
Under $H_{\Fact}(X)$, the uses of $\Fact(X)$ in
Theorem~\ref{thm:A-infinity-2}
are stable, presentable, symmetric monoidal, and compatible with
colimits in each variable.
\end{proposition}

\begin{proof}
The de~Rham sheaf ambient is stable and presentable in the GR
formalism \cite[Chapter~I.3]{GR17}; Francis--Gaitsgory construct the
chiral monoidal operation on $D(\Ran X)$ \cite{Francis2012}.  Clauses
(b)--(d) specify the model presentation used in this volume.  The
properadic transfer adds the separate package $H_{\mathrm{prop}}$ of
Theorem~\ref{thm:factorization-properad-transfer}.
\end{proof}

%================================================================
\section{Simplicial properads and factorization transfer}
\label{sec:ainf2-properads}
%================================================================

An operad records many inputs and one output.  A properad records many
inputs and many outputs, with composition indexed by connected graphs
\cite{Val07}.  We use the profile convention
$\cP(n,m)=$ ``$n$ inputs and $m$ outputs.''

A Verdier--Koszul properadic datum consists of a specified evaluation
and its duality transpose
\[
 \operatorname{ev}_{A}\colon
 A\otimes^{\ch}K_X(A)\longrightarrow\omega_X,
 \qquad (2,1),
\]
\[
 \operatorname{coev}_{A}\colon
 \omega_X\longrightarrow K_X(A)\otimes^{\ch}A,
 \qquad (1,2),
\]
where $K_X(A)=\mathbb D_{\Ran}B_X(A)$.  The zig--zag identities are
part of the dual-pair clause of $H_{\mathrm{VD}}$.  A multiplication
or Maurer--Cartan operation has profile $(n,1)$; its Verdier transpose
has profile $(1,n)$.  These profiles generate the properadic envelope
used below.

\begin{definition}[Factorization-enriched properad;
\ClaimStatusDefinitional]
\label{def:v1-factorization-properad}
\index{properad!factorization}
Let $\cV_X$ be the symmetric-monoidal model enhancement selected by
$H_{\Fact}(X)$.  A \emph{$\cV_X$-enriched properad} is a colored
collection $\{\cP(\mathbf d;\mathbf c)\}$ in $\cV_X$, equivariant in
the input and output colors, with units and associative substitution
maps indexed by connected directed graphs.  It is a
\emph{factorization-enriched properad} when every profile satisfies
disjoint-locus factorization descent and the substitution maps commute
with those descent equivalences.  We write $\FactProp(X)$ for the
resulting category whenever the model presentation of
Theorem~\ref{thm:factorization-properad-transfer} is available.
\end{definition}

\begin{theorem}[Model-categorical presentation of simplicial
properads; \ClaimStatusConjectured]
\label{thm:hackney-robertson-model}
\textup{[Type signature: external homotopy-theory input; simplicial
properad presentation; properadic enrichment level; hypothesis
package: the presentation clause $H_{\mathrm{HR}}$ stated in the
theorem body.]}
The presentation clause $H_{\mathrm{HR}}$ asserts: the ordinary
category $\mathsf{sProperad}$ of small simplicial properads carries a
cofibrantly generated model structure in which a morphism
$f\colon P\to Q$ is a weak equivalence precisely when it is a weak
homotopy equivalence on every input--output profile and
$\pi_0f\colon\pi_0P\to\pi_0Q$ is an equivalence of categories, and a
fibration precisely when it is a Kan fibration on every profile and
$\pi_0f$ is an isofibration; and this model structure presents the
homotopy theory of $\infty$-properads of
Hackney--Robertson--Yau \cite{HRY17}.
\end{theorem}

\begin{remark}[Scope of the literature input]
\label{rem:hackney-robertson-scope}
Hackney--Robertson--Yau \cite{HRY17} construct the homotopy theory of
$\infty$-properads inside graphical sets, a presheaf model built from
simplicial sets; that construction does not itself furnish the model
structure on $\mathsf{sProperad}$ displayed in $H_{\mathrm{HR}}$.
The clause $H_{\mathrm{HR}}$ is therefore consumed as a named
hypothesis, with the graphical-set homotopy theory as its evidence,
and every use below routes through
Theorem~\ref{thm:factorization-properad-transfer}, whose package
$H_{\mathrm{prop}}$ is stated independently.
\end{remark}

\begin{theorem}[Conditional factorization-enriched properad transfer;
\ClaimStatusConditional]
\label{thm:factorization-properad-transfer}
\textup{[Type signature: Open quadrant; factorization-enriched
properad presentation; Beilinson level $1\leftrightarrow2$ envelope;
hypothesis package $H_{\mathrm{prop}}$.]}
Let $\cV_X$ be a cofibrantly generated symmetric-monoidal model
enhancement of the factorization ambient.  The package
$H_{\mathrm{prop}}$ consists of the following clauses.
\begin{enumerate}[label=\textup{(P\arabic*)}]
\item \label{Hprop-monoid} $\cV_X$ satisfies the monoid axiom.
\item \label{Hprop-unit} The monoidal unit of $\cV_X$ is cofibrant.
\item \label{Hprop-admissible} The free properad monad formed by
connected-graph substitution is admissible: its generating cell
domains are small and the transferred factorizations exist.
\item \label{Hprop-weq} Connected-graph extension and the operad-to-
properad transfer preserve weak equivalences between cofibrant
collections; on the envelope generated by the associative bar--cobar
pair they preserve the derived unit and counit equivalences.
\item \label{Hprop-descent} The localization imposing disjoint-locus
factorization descent is stable under the free properad monad,
homotopy colimits of cofibrant diagrams, and the transferred
factorizations.
\end{enumerate}
Under $H_{\mathrm{prop}}$, factorization-enriched properads carry the
transferred cofibrantly generated model structure.  Its weak
equivalences are the profilewise weak equivalences together with the
equivalence condition on the color category.  The graphwise extension
of the associative bar--cobar adjunction is a Quillen equivalence on
the properadic envelope generated by that pair.
\end{theorem}

\begin{proof}
Clauses~\ref{Hprop-monoid}--\ref{Hprop-admissible} transfer the model
structure from enriched collections to algebras over the free
properad monad.  Clause~\ref{Hprop-descent} descends this structure to
the factorization-local objects.  Clause~\ref{Hprop-weq} carries the
derived bar--cobar unit and counit through connected-graph extension,
which gives the asserted Quillen equivalence on the generated
envelope.  The presentation clause $H_{\mathrm{HR}}$ of
Theorem~\ref{thm:hackney-robertson-model} posits the
simplicial properad model that motivates the profile and color
conditions; the enrichment over $\cV_X$ is governed by the five
displayed clauses.
\end{proof}

%================================================================
\section{Conditional factorization and properadic forms of Theorem~A}
\label{sec:ainf2-main}
%================================================================

\begin{theorem}[Factorization restriction and properadic envelope of
Theorem~A; \ClaimStatusConditional]
\label{thm:A-infinity-2}
\label{thm:theorem-A-E1}
\index{Theorem A-infinity-2}
\index{Theorem A!E_1-ordered variant}
\textup{[Type signature: Open quadrant; Ran, factorization, and
properadic presentations; Beilinson level $1\leftrightarrow2$;
pro-nilpotent chiral Ran ambient; hypothesis packages
$H_{\Fact}+H_{\mathrm{fact}}+H_{\mathrm{conv}}$, with
$H_{\mathrm{prop}}$ for clause~\ref{A2-i}.]}
Let $X$ be a smooth curve over a characteristic-zero field and place
the reduced associative operad in
$\mathcal C_X=(D(\Ran X),\otimes^{\ch})$.  Assume
$H_{\Fact}(X)$, $H_{\mathrm{fact}}$, and $H_{\mathrm{conv}}$.
Enhanced bar and cobar restrict to a derived
adjoint equivalence
\[
 \Omegach_X\colon\CoFact^{\conil,\comp}(X)
 \rightleftarrows
 \Fact^{\aug,\comp}(X)\colon\Bbarch_X
\]
on the chosen factorization subcategories, with unit and counit
equivalences.  It has the following refinements.
\begin{enumerate}[label=\textup{(\roman*)}]
\item \label{A2-i} \emph{Properadic envelope;
$\ClaimStatusConditional$.}
Under $H_{\mathrm{prop}}$, connected-graph extension gives a derived
adjoint equivalence
\[
 \Omegach_X\colon
 \CoFactProp^{\conil,\comp}_{\langle\chirAss\rangle}(X)
 \rightleftarrows
 \FactProp^{\aug,\comp}_{\langle\chirAss\rangle}(X)
 \colon\Bbarch_X .
\]
The subscript $\langle\chirAss\rangle$ denotes the envelope generated
from the associative operadic pair by connected-graph substitution,
radical completion, and weak equivalence.

\item \label{A2-ii} \emph{Fiber restriction;
$\ClaimStatusConditional$.}
For a marked point $x\in X$, assume that the fiber functor
$i_x^!\colon\Fact(X)\to\mathsf{Ch}(k)$ is exact symmetric monoidal and
commutes with the completed bar and cobar constructions.  Its image is
the ordinary associative bar--cobar adjunction, whose unit and counit
are quasi-isomorphisms \cite[Corollary~2.3.4]{LV12}.

\item \label{A2-iii} \emph{Ordered-to-symmetric descent;
$\ClaimStatusConditional$.}
Let $R(z)$ satisfy hypotheses~\ref{R1} and~\ref{R2} of
Lemma~\ref{lem:R-twisted-descent-unitary} and arise as monodromy of the
stated flat connection.  On the open configuration locus the principal
$\Sigma_n$-torsor
\[
 \pi_n^\circ\colon\Conf_n^{\ord}(X)\longrightarrow\Conf_n(X)
\]
gives $L_R$-twisted descent of the ordered bar.  Under
$H_{\mathrm{fact}}^{R}$, nearby cycles extend this equivalence through
the Fulton--MacPherson collision strata to the cardinality-$n$ Ran
stratum.
\end{enumerate}
\end{theorem}

\begin{proof}
Francis--Gaitsgory's Proposition~4.1.2, together with the
pro-nilpotence argument in the proof of their Theorem~5.1.1, gives
the enhanced equivalence in $D(\Ran X)$.
The closure maps of $H_{\mathrm{fact}}$ restrict it to the displayed
factorization subcategories.

On chain models the two free constructions are
\[
 \Bbarch_X(A)=\Tc_{\Fact(X)}(s^{-1}\bar A),
 \qquad
 \Omegach_X(C)=T_{\Fact(X)}(s\bar C).
\]
The bar multiplication term and the cobar coproduct term both have
cohomological degree~$+1$ by
Convention~\ref{conv:A-infinity-2-ordered-bar-signs}.  Their universal
twisting-morphism property gives the adjunction
\cite[Theorem~2.2.9]{LV12}.

The three convergence clauses enter at distinct points.  The direct-
sum bar is the filtered colimit of its finite length pieces by
$H_{\mathrm{len}}$.  The identity
$A\simeq\varprojlim_qA/\bar A^{\,q}$ is supplied by
$H_{\mathrm{aug}}$ and identifies the completed algebra with its
finite augmentation quotients.  The pro-coradical tower of
$H_{\mathrm{ML}}$ has vanishing derived inverse-limit obstruction, so
the stagewise unit and counit equivalences pass to the completed limit.
This proves the first assertion.

Clause~\ref{A2-i} is Theorem~\ref{thm:factorization-properad-transfer}.
The exact symmetric-monoidal fiber functor in clause~\ref{A2-ii}
carries the displayed free constructions to $T^c(s^{-1}\bar A)$ and
$T(s\bar C)$; Loday--Vallette's unit and counit theorem then applies.
Clause~\ref{A2-iii} is the open-locus assertion and the
$H_{\mathrm{fact}}^{R}$ extension of
Lemma~\ref{lem:R-twisted-descent-unitary}.
\end{proof}

\begin{remark}[Three convergence mechanisms]
\label{rem:A-inf-2-ml-step-2}
\index{Mittag-Leffler!bar-length vs coradical filtration}
For the direct-sum bar,
\[
 F^{\mathrm{bar}}_{\leq p}B(A)
 =\bigoplus_{0\leq n\leq p}(s^{-1}\bar A)^{\otimes n},
 \qquad
 B(A)=\varinjlim_pF^{\mathrm{bar}}_{\leq p}B(A).
\]
This is bar-length exhaustiveness.  Augmentation completeness is the
separate comparison
\[
 A\xrightarrow{\sim}\varprojlim_qA/\bar A^{\,q}.
\]
Coradical convergence belongs to a pro-coalgebra presentation
$C\simeq\varprojlim_qC_{\leq q}$.  For every cohomological degree and
every properadic profile, $H_{\mathrm{ML}}$ requires the tower
$\{H^j(C_{\leq q})\}_q$ to be Mittag--Leffler.  The Milnor sequence
\[
 0\longrightarrow
 R^1\!\varprojlim_q H^{j-1}(C_{\leq q})
 \longrightarrow H^j(R\!\varprojlim_q C_{\leq q})
 \longrightarrow \varprojlim_q H^j(C_{\leq q})
 \longrightarrow0
\]
then identifies the derived inverse limit with the ordinary limit.
Thus $H_{\mathrm{len}}$, $H_{\mathrm{aug}}$, and $H_{\mathrm{ML}}$
control respectively a filtered colimit, an algebraic completion, and
a derived inverse limit.
\end{remark}

\begin{remark}[Fixed curve, relative interior, and Deligne--Mumford boundary]
\label{rem:A-infinity-2-modular-family-scope}
Theorem~\ref{thm:A-infinity-2} is a fixed-smooth-curve theorem.  It
has two relative enhancements.  The enhancement over
$\mathcal{M}_{g,n}$ carries base change for the relative Ran prestack;
the enhancement over $\overline{\mathcal{M}}_{g,n}$ carries, in
addition, nodal factorization sewing.

Let $\pi:\mathcal C\to S$ be a smooth projective family of marked
curves and let $\Ran^{\rel}(\mathcal C/S)\to S$ be the relative Ran
prestack. For $f:S'\to S$, write
$f':\Ran^{\rel}(\mathcal C'/S')\to\Ran^{\rel}(\mathcal C/S)$ for the
induced map. The relative interior obstruction is the base-change
comparison
\[
\alpha_{\mathrm{BC}}:\;
f^{!}\pi^{\Ran}_{*}\Longrightarrow \pi'^{\Ran}_{*}f'^{!}
\]
on relative factorization sheaves.  For
$\mathcal F\in\Fact^{\rel}(\mathcal C/S)$ satisfying
$H_{\mathrm{fact}}+H_{\mathrm{conv}}$, the cone of this comparison
defines
\[
[\alpha_{\mathrm{BC}}]\in
\mathrm{Ext}^{1}_{\Shv^{\DR}(S')}
\bigl(f^{!}\pi^{\Ran}_{*}\mathcal F,\;
\pi'^{\Ran}_{*}f'^{!}\mathcal F\bigr).
\]
Vanishing of this class for the conilpotent-complete subcategory gives
the relative version of Theorem~\ref{thm:A-infinity-2} over
$\mathcal M_{g,n}$.

Let $D_\Delta\subset\overline{\mathcal M}_{g,n}$ be a boundary divisor
with normalisation $\widetilde C=C'\sqcup C''$ and formal node
$\widehat D_p$. Mok's logarithmic Fulton--MacPherson geometry
\cite[Theorem~3.3.1, Theorem~5.2.3, Corollary~5.3.4]{Mok25} supplies
the normal-crossings and degeneration space on which nodal sewing is
formulated. The chain-level factorization statement is the separate
nearby-cycle comparison
\[
\beta_{\mathrm{nodal}}:\;
\Psi_\Delta(\pi_{\log})_{*}\mathcal F
\Longrightarrow
(\pi_{\log})_{*,D_\Delta}
\bigl(\mathcal F(C')\otimes_{\mathcal F(\widehat D_p)}
\mathcal F(C'')\bigr),
\]
where $\Psi_\Delta$ denotes Beilinson--Drinfeld factorization nearby
cycles. Its obstruction class is
\[
[\beta_{\mathrm{nodal}}]\in
\mathrm{Ext}^{1}_{\Shv^{\DR}(D_\Delta)}
\bigl(\Psi_\Delta(\pi_{\log})_{*}\mathcal F,\;
(\pi_{\log})_{*,D_\Delta}
\bigl(\mathcal F(C')\otimes_{\mathcal F(\widehat D_p)}
\mathcal F(C'')\bigr)\bigr).
\]
Vanishing of $[\beta_{\mathrm{nodal}}]$ at every boundary divisor is
the extra condition for the Deligne--Mumford extension.

The two obstructions are independent. The class
$[\alpha_{\mathrm{BC}}]$ controls horizontal transport over smooth
families; $[\beta_{\mathrm{nodal}}]$ controls sewing at nodal
degenerations. Class~$\mathsf G$ satisfies the nodal condition by
direct free-field sewing. Class~$\mathsf L$ satisfies it under the
generic non-critical Beilinson--Drinfeld chiral sewing hypotheses.
For class~$\mathsf M$, the sewing datum lives in the weight-completed
lane.

The $R$-twisted descent of Lemma~\ref{lem:R-twisted-descent-unitary} is a
fixed-curve configuration-space theorem.  Its Fulton--MacPherson
collision extension is the clause $H_{\mathrm{fact}}^{R}$.  A
Deligne--Mumford extension additionally carries the family classes
$[\alpha_{\mathrm{BC}}]$ and $[\beta_{\mathrm{nodal}}]$.  Their
vanishing extends the family comparison across the compactified base.
Together with the indicated Theorem~D packages, the resulting hierarchy
is
\[
 \operatorname{Obs}^{\mathrm{def}}_g(\cA)
 \in H^2\!\left(\Def_g(\cA)\right),
 \qquad g\geq2,
 \qquad
 \operatorname{tr}_1\operatorname{Obs}^{\mathrm{def}}_{1,1}(\cA)
 =+\kappa(\cA)\lambda_1
 \in H^2(\overline{\cM}_{1,1})
 \quad\text{under }H_D^1,
\]
\[
 \mathfrak O_g^K(\cA)
 =\kappa(\cA)\lambda_{-1}(\mathbb E_g)
 \in K^0(\overline{\cM}_g),
 \qquad
 \operatorname{ch}_g\!\left(\mathfrak O_g^K(\cA)\right)
 =(-1)^g\kappa(\cA)\lambda_g
 \in H^{2g}(\overline{\cM}_g)
 \quad\text{under }H_D^K.
\]
The package~$H_D^{\mathrm{tr}}$ supplies the derived comparison between
the deformation and virtual Hodge carriers.  Under
$H_D^{\mathrm{graph}}$, the stable-graph scalar is
\[
 F_g(\cA)=\kappa(\cA)\lambda_g^{\mathrm{FP}}
 +\delta F_g^{\mathrm{cross}}(\cA)\in\mathbb C.
\]
\end{remark}

\begin{remark}[Alternative proof lanes: Theorem~A versus Theorem~C]
\label{rem:A-inf-2-ptvv-scope}
\index{Theorem A!alternative proof lanes}
\index{PTVV!Theorem C proof lane}
Lurie~HA~\S5.5 gives a topological route to the locally constant
specialization of clause~\ref{A2-ii}: factorization algebras via little
discs and factorization envelopes reproduce the classical associative
Koszul pair.  The $\PTVV$ shifted-symplectic construction belongs to
\emph{Theorem~C}: its lane
constructs the $(-(3g-3))$-shifted symplectic structure on
$\RMap(\overline{\cM}_g,BG_\cA)$ as a genuinely disjoint path from the
fiber-center Koszul identification, relevant to the characteristic-datum
C0/C1/C2 clauses and to the $\lambda_g$-free formulation of Theorem~D.
Francis--Gaitsgory and Loday--Vallette supply the bar--cobar
reconstruction inputs.  The simplicial properad model structure is the
presentation clause $H_{\mathrm{HR}}$
(Theorem~\ref{thm:hackney-robertson-model}, Conjectured), while
$H_{\mathrm{prop}}$ carries its factorization-enriched transfer.
\end{remark}

\begin{corollary}[Classical chain model from fiber restriction;
\ClaimStatusConditional]
\label{cor:classical-A-from-A-infinity-2}
\index{Theorem A!descent from A-infinity-2}
Assume the fiber and Verdier clauses of
Theorem~\ref{thm:A-infinity-2}\ref{A2-ii}, $H_{\mathrm{VD}}$, and the
following fiber package $H_{\mathrm{fib}}$: a chosen !-fiber functor
on the finite-window holonomic subcategory is exact, its tensor
exchange maps are isomorphisms, it commutes with the completed inverse
limits used by bar and cobar, and Verdier base change identifies the
fiber of $\mathbb D_{\Ran}$ with the linear dual of the fiber.  The
geometric bar--cobar adjunction
\textup{(}Theorem~\textup{\ref{thm:bar-cobar-isomorphism-main}})
and its Verdier-duality intertwining
\textup{(}Theorem~\textup{\ref{thm:bar-cobar-verdier}}) are the images
of the enhanced Ran constructions under this fiber functor.
\end{corollary}

\begin{proof}
Clause~\ref{A2-ii} and $H_{\mathrm{fib}}$ carry $\Bbarch_X$ and $\Omegach_X$ to
$T^c(s^{-1}\bar A)$ and $T(s\bar C)$ and carries their unit and counit
to the classical maps.  The base-change, K\"unneth, and biduality maps
of $H_{\mathrm{VD}}$ carry the Verdier comparison to the same fiber.
\end{proof}

%================================================================
\section{$R$-matrix-twisted $\Sigma_n$-descent}\label{sec:ainf2-R-descent}
%================================================================

The descent lemma relating ordered and symmetric bar for
$\Eone$-chiral algebras is the explicit content of
Corollary~\ref{cor:eight-cor-R-descent} of the unified Koszul
reflection (equivalently, of
Theorem~\ref{thm:A-infinity-2}\ref{A2-iii}).

\begin{proposition}[Classical MC/CYBE output and quantum QYBE input;
\ClaimStatusConditional]
\label{prop:classical-cybe-quantum-r-input}
\index{Maurer-Cartan!classical output versus quantum input}
\index{CYBE!classical output}
\index{QYBE!quantum input}
The universal twisting morphism of
Corollary~\ref{cor:eight-cor-twisting-morphism} produces a
classical Maurer--Cartan equation in the ordered convolution algebra,
\[
D\theta_\cA+\tfrac12[\theta_\cA,\theta_\cA]=0.
\]
On the binary collision stratum, if the residue kernel is
$r(z)\in(\End \cA^{\otimes 2})((z))$, the flatness of the classical
logarithmic connection
\[
\nabla_{\mathrm{cl}}
=d-\hbar\sum_{i<j}r_{ij}(z_i-z_j)\,d(z_i-z_j)
\]
is the spectral classical Yang--Baxter identity
\[
[r_{12}(z_{12}),r_{13}(z_{13})]
+[r_{12}(z_{12}),r_{23}(z_{23})]
+[r_{13}(z_{13}),r_{23}(z_{23})]=0
\]
together with the Arnold relation for the logarithmic forms when
$r(z)$ has rational KZ form $\Omega/z$. This is classical output of
the bar differential. A quantum descent datum is extra input: a
formal or meromorphic family
\[
R(z)=1+\hbar r(z)+O(\hbar^2),\qquad
\check R(z):=\tau\circ R(z),
\]
whose checked operators satisfy the braid cocycle and whose ordinary
operators satisfy the quantum Yang--Baxter equation and unitarity of
Lemma~\ref{lem:R-twisted-descent-unitary}.  The $\hbar^2$ coefficient
of QYBE is CYBE.  The coefficients $R^{(m)}(z)$ for $m\geq2$ and the
flat realization of $L_R$ form the additional quantum descent datum.
\end{proposition}

\begin{proof}
Expanding QYBE for
$R(z)=1+\hbar r(z)+\hbar^2R^{(2)}(z)+O(\hbar^3)$ gives identical
linear terms on both sides. At order~$\hbar^2$, the unknown
$R^{(2)}$ terms cancel, and the remaining difference is precisely
the displayed sum of commutators. Expanding unitarity
$R^{21}(-z)R^{12}(z)=1$ at order~$\hbar$ gives
$r^{21}(-z)+r^{12}(z)=0$. The Maurer--Cartan equation is the same
flatness identity before choosing a quantization; the Arnold relation
annihilates the three logarithmic two-form coefficients on
$\Conf_3(X)$.  Thus the expansion determines the classical kernel,
while the higher coefficients $R^{(m)}(z)$ and their flat connection
are quantum input.
\end{proof}

\begin{lemma}[$R$-matrix-twisted $\Sigma_n$-descent, unitary quantum $R$;
\ClaimStatusConditional]
\label{lem:R-twisted-descent-unitary}% renamed 2026-05-17 from lem:R-twisted-descent (collision with ordered_associative_chiral_kd.tex:587)
\index{R-matrix!Sigma-descent}
\index{descent!R-twisted}
\textup{[Type signature: Open quadrant; ordered/symmetric
configuration presentation; Beilinson level~$2$; hypotheses
\ref{R1}+\ref{R2}, with $H_{\mathrm{fact}}^{R}$ for Ran extension.]}
Let $\cA$ be an $\Eone$-chiral algebra with a meromorphic unitary
quantum $R$-matrix
$R(z) \in \operatorname{Aut}(\cA \otimes \cA)((z))$; write
$\check R(z):=\tau\circ R(z)$ for the checked braiding. Assume
\begin{enumerate}[label=\textup{(R\arabic*)}]
\item \label{R1} the quantum Yang--Baxter equation
$R_{12}(z_{12})\,R_{13}(z_{13})\,R_{23}(z_{23})
=R_{23}(z_{23})\,R_{13}(z_{13})\,R_{12}(z_{12})$ on
$\{z_1,z_2,z_3\}$-triples; and
\item \label{R2} \emph{unitarity}: $R^{21}(-z)\,R^{12}(z)=\id_{\cA\otimes\cA}$
on the formal punctured disc, where
$R^{21}(w):=\tau\cdot R(w)\cdot\tau$ is the swap-conjugate.
\end{enumerate}
Assume also that the meromorphic family is realized as the monodromy
of a flat KZ/Cherednik-type connection on
$\Conf^{\ord}_n(X)$ for each~$n$, with elementary chamber-crossing
operator $\check R_i(z_i-z_{i+1})$ on adjacent factors.
Then
\[
 \pi_n^\circ\colon\Conf_n^{\ord}(X)\longrightarrow\Conf_n(X)
\]
is a principal $\Sigma_n$-torsor, and there is a canonical
$\Sigma_n$-equivariant local system $L_R$ on
$\Conf^{\ord}_n(X)$ whose descent isomorphism for the elementary
transposition $s_i$ is $\check R_i(z_i-z_{i+1})$.  On the open locus,
\[
 \left.\barB^{\Sigma}_n(\cA)\right|_{\Conf_n(X)}
 \simeq
 \operatorname{Desc}_{\pi_n^\circ}
 \left(\left.B^{\ord}_n(\cA)\right|_{\Conf_n^{\ord}(X)}\otimes L_R\right).
\]
Under $H_{\mathrm{fact}}^{R}$, $L_R$ has compatible nearby cycles on
the Fulton--MacPherson collision strata, and the open equivalence
extends to
\[
\barB^{\Sigma}_n(\cA)
\;\simeq\;
\bigl(B^{\ord}_n(\cA)\bigr)^{\Sigma_n\text{-coinv},\; L_R\text{-twisted}}.
\]
\end{lemma}

\begin{proof}
The free action of $\Sigma_n$ on
$\Conf_n^{\ord}(X)=X^n\setminus\Delta$ permutes the labels, and its
geometric quotient is $\Conf_n(X)$.  Hence $\pi_n^\circ$ is a
principal $\Sigma_n$-torsor.

Assign to each point
$(z_1,\dots,z_n) \in \Conf^{\ord}_n(X)$ the vector space
$\cA^{\otimes n}$ at that point and to each elementary deck
transformation $s_i$ the isomorphism
\[
\phi_i=\check R_i(z_i-z_{i+1})\colon
s_i^*(\cA^{\otimes n})\longrightarrow \cA^{\otimes n}.
\]
The QYBE of~\ref{R1} is exactly the braid cocycle
\[
\phi_i\,s_i^*(\phi_{i+1})\,(s_i s_{i+1})^*(\phi_i)
=
\phi_{i+1}\,s_{i+1}^*(\phi_i)\,(s_{i+1}s_i)^*(\phi_{i+1})
\]
on triple overlaps; distant generators commute because they act on
disjoint tensor factors. The unitarity condition~\ref{R2} gives
$\phi_i\,s_i^*(\phi_i)=\id$, since $s_i$ sends $z_i-z_{i+1}$ to
$-(z_i-z_{i+1})$ and swaps the two factors. Thus
\ref{R1}+\ref{R2} supply the Coxeter descent cocycle on the ordered
cover. The flat KZ/Cherednik realization supplies the underlying
local system on $\Conf^{\ord}_n(X)$; the cocycle above supplies its
$\Sigma_n$-equivariant descent datum.

Descent along the principal torsor is the standard equivalence between
sheaves on $\Conf_n(X)$ and sheaves on the ordered cover equipped
with a $\Sigma_n$-cocycle. The ordered bar $B^{\ord}_n(\cA)$ restricted to
$\Conf^{\ord}_n(X)$ is
$\cA^{\boxtimes n} \otimes \Omega^{n-1}_{\Conf^{\ord}_n(X)}(\log D)$,
and its descent to $\Conf_n(X)$ with twist by $L_R$ is by definition
the symmetric bar $\barB^{\Sigma}_n(\cA)$ of
Definition~\ref{def:geometric-bar} (the bar coefficient in the
symmetric setting is the Arnold 1-form $d\log(z_i - z_j)$, which
transforms under $\Sigma_n$ by the sign character twisted by
the checked operator $\check R(z_i-z_j)$ acting on the fibres).

The log Fulton--MacPherson compactification
$\overline{\Conf}_n(X)$ supplies the fixed-curve normal-crossings
boundary on which $H_{\mathrm{fact}}^{R}$ is formulated.  That clause
identifies the nearby-cycle monodromy with the chiral collision
product.  The ordered bar differential is a coderivation for
deconcatenation and preserves stratum depth, so descent on a boundary
stratum is the top-stratum descent applied to the nearby-cycle fibre.
For the pole-free case the exceptional-divisor limit is the checked
swap $\tau$, and
the formula reduces to ordinary $\Sigma_n$-coinvariants.
\end{proof}

\begin{remark}[Unitarity and collision extension]
\label{rem:R-descent-unitarity-scope}
\index{R-matrix!unitarity scope}
\index{descent!scope across families}
Hypothesis~\ref{R2} is satisfied on the finite-dimensional evaluation
module lane of the rational Drinfeld Yangian $Y_\hbar(\fg)$ after the
standard scalar normalization of the rational $R$-matrix, and on the
integrable finite-type lane of trigonometric quantum affine
$U_q(\hat\fg)$ at generic~$q$ after the corresponding normalization.
These are representation-surface statements.  A factorization
Drinfeld--Kohno equivalence on Yangian category~$\mathcal O$ and a
dg-shifted Yangian/line-operator comparison carry their own comparison
morphisms.  At roots of unity, and for elliptic or dynamical
$R$-matrices of Belavin/Felder type, the chosen presentation supplies
unitarity and the regular-singular boundary extension.  For chiral
algebras in a genuinely trivial ordinary
$R$-matrix lane~$R(z)=1$, the checked braiding is $\check R(z)=\tau$;
unitarity~\ref{R2} is immediate and~\ref{R1} is vacuous. The
non-zero-level Heisenberg lane is separate: its ordered-bar braiding is
the scalar $R_{\cH_k}(z)=\exp(k\hbar/z)$. It satisfies
unitarity~\ref{R2} by scalar cancellation and satisfies~\ref{R1}
tautologically; its local system becomes trivial on the zero-level and
zero-charge specializations.

The fixed-curve collision extension is $H_{\mathrm{fact}}^{R}$.  For
the rational KZ and trigonometric Cherednik connections it is the
regular-singular logarithmic extension along the Fulton--MacPherson
boundary on the stated representation lanes.  Elliptic and dynamical
$R$-matrices use the chosen compactification and dynamical-parameter
convention.  Thus hypotheses~\ref{R1}+\ref{R2} govern open
configuration-space descent, while $H_{\mathrm{fact}}^{R}$ governs
the Ran-level symmetric bar.
\end{remark}

\begin{remark}[Comparison with $\fg^{\Eone}_\cA$ versus $\fg^{\mathrm{mod}}_\cA$]
\label{rem:av-via-descent}
\index{averaging map!via R-twisted descent}
The averaging map
$\mathrm{av}\colon \fg^{\Eone}_\cA \to \fg^{\mathrm{mod}}_\cA$
(\S\ref{sec:e1-first-prose-architecture}) is the dual of the descent
map $\pi^!$ after forgetting the monodromy datum: dualising
$\pi\colon \Ran^{\ord}(X) \to \Ran(X)$ and tensoring with the structure
sheaf gives the $\Sigma_n$-coinvariant projection on the classical
shadow.  Quantization is supplied by the quantum $R$-datum, and the
full modular characteristic additionally uses the conformal-vector
normalization. For affine
Kac--Moody $V_k(\fg)$ in trace-form convention, the classical residue
kernel is
\[
r^{KM}(z)=k\,\Omega_{\mathrm{tr}}/z.
\]
Its trace-form contribution to the scalar diagonal is
$k\,\dim(\fg)/(2h^\vee)$. The full invariant is
\[
\kappa(V_k(\fg))=\dim(\fg)(k+h^\vee)/(2h^\vee),
\]
where the trace of the descent local system contributes
$k\dim(\fg)/(2h^\vee)$ and the Sugawara conformal vector contributes
$\dim(\fg)/2$.
\end{remark}

%================================================================
\section{Typed restrictions of Theorem~A}\label{sec:ainf2-downstream}
%================================================================

The Ran equivalence, its factorization restriction, its completed
chain model, its conditional properadic transfer, and its ordered
descent occupy five distinct hypothesis packages.

\begin{corollary}[Five typed outputs of Theorem~A;
\ClaimStatusConditional]
\label{cor:ainf2-downstream-list}
\index{Theorem A-infinity-2!downstream corollaries}
The following outputs hold on their stated domains:
\begin{enumerate}[label=\textup{(D\arabic*)}]
\item \label{D1} In the pro-nilpotent chiral Ran ambient, enhanced
associative bar and cobar are inverse equivalences in $D(\Ran X)$
\textup{(}Theorem~\ref{thm:koszul-reflection}\ref{KR-i}\textup{)}.
\item \label{D2} Under $H_{\mathrm{fact}}$, the equivalence restricts
to the chosen associative factorization subcategories
\textup{(}Theorem~\ref{thm:koszul-reflection}\ref{KR-iii}\textup{)}.
\item \label{D3} Under $H_{\mathrm{conv}}$, the completed chain models
have bar $T^c(s^{-1}\bar A)$, cobar $T(s\bar C)$, and unit and counit
equivalences \textup{(}Theorem~\ref{thm:A-infinity-2}\textup{)}.
\item \label{D4} Under $H_{\mathrm{prop}}$, connected-graph extension
gives the Quillen equivalence on the properadic envelope
\textup{(}Theorem~\ref{thm:A-infinity-2}\ref{A2-i}\textup{)}.
\item \label{D5} Under hypotheses~\ref{R1}+\ref{R2}, the ordered bar
descends along the principal $\Sigma_n$-torsor on the open
configuration locus; $H_{\mathrm{fact}}^{R}$ supplies its Ran
extension \textup{(}Theorem~\ref{thm:A-infinity-2}\ref{A2-iii}\textup{)}.
\end{enumerate}
\end{corollary}

\begin{proof}
Items~\ref{D1}--\ref{D3} are the Ran equivalence, its factorization
closure, and the convergence argument of
Theorem~\ref{thm:A-infinity-2}.  Item~\ref{D4} is
Theorem~\ref{thm:factorization-properad-transfer}; item~\ref{D5} is
Lemma~\ref{lem:R-twisted-descent-unitary}.
\end{proof}

\begin{remark}[Cross-volume bridges to climax and conductor chapters]
\label{rem:cross-volume-bridges-to-climax-and-conductor}
\index{cross-volume bridge!climax and conductor}
The unified Koszul reflection
(Theorem~\ref{thm:koszul-reflection}) feeds directly into two
parallel-established chapters:
\begin{itemize}
\item \textbf{Climax chapter} \textup{(}\ref{chap:climax-platonic})
identifying $\bar\partial=\KZ^{*}(\nabla_{\Arnold})$ as the bar
differential on the universal Arnold connection: the unit~$\eta$ of
the Koszul reflection at the geometric Arnold connection produces the
Knizhnik--Zamolodchikov pullback identification of the bar
differential.
\item \textbf{Universal conductor chapter}
\textup{(}\ref{chap:universal-conductor}) proving
$K=-c_{\mathrm{ghost}}(\mathrm{BRST})$ as the universal ghost-charge
conductor: the conductor
$K^{\kappa}(\cA)=\kappa(\cA)+\kappa(\cA^{!})$ is the numerical
observable of the Verdier-dualized Koszul pair.  The functor
$K_X=\mathbb D_{\Ran}B_X$ is its categorical antecedent.  Pairwise
equivalences
$K_X(\cA)\simeq\cA^!$ and $K_X(\cA^!)\simeq\cA$
belong to the separate represented self-duality package.
\end{itemize}
\end{remark}

%================================================================
\section{Symmetric monoidal ambient corrections}\label{sec:ainf2-heals}
%================================================================

Theorem~\ref{thm:A-infinity-2} uses the following ambient data.

\begin{remark}[Symmetric-bar and ordered-bar forms]
\label{rem:ainf2-heal-69}

The factorization form has two parts:
\begin{itemize}
\item[\textsc{Symmetric}:] The \emph{symmetric-bar} statement lies
in the pro-nilpotent Francis--Gaitsgory Ran ambient; its
factorization restriction uses $H_{\mathrm{fact}}$.
\item[\textsc{Ordered}:] The \emph{ordered-bar} statement is
\ref{A2-iii} of Theorem~\ref{thm:A-infinity-2}: the
ordered bar on $\Conf_n^{\ord}(X)$ carries the $R$-twisted local system
$L_R$; descent along $\pi_n^\circ$ gives the symmetric form on
$\Conf_n(X)$, and $H_{\mathrm{fact}}^{R}$ gives the Ran extension.
\end{itemize}
Both forms hold on their stated hypotheses. The level-$0$
primitive is the open factorization category and its chosen vacuum;
the level-$2$ ordered object is $B^{\ord}(\cA)$ on $\Ran^{\ord}(X)$,
and the symmetric bar is the $L_R$-twisted $\Sigma_n$-coinvariant
descent image.
\end{remark}

\begin{remark}[Factorization unit functor]
\label{rem:ainf2-heal-70}

The operadic transfer uses the symmetric-monoidal functor
\[
 \mathsf{Ch}(k)\longrightarrow\Fact(X),
 \qquad M\longmapsto M\otimes\mathbf 1_{\Fact(X)}.
\]
Here $\mathbf 1_{\Fact(X)}$ is the factorization unit represented by
the constant factorization algebra on $\cO_X$.  This is the unit
functor in $H_{\Fact}(X)$.
\end{remark}

\begin{remark}[$R$-twisted descent]
\label{rem:ainf2-heal-71}

For $\cH_k$ the ordered bar carries the scalar descent datum
\[
R_{\cH_k}(z)=\exp(k\hbar/z),
\]
and Lemma~\ref{lem:R-twisted-descent-unitary} identifies its symmetric
bar with the scalar-local-system descent.  The zero-level and
zero-charge specializations give the trivial local system.
\end{remark}

\begin{remark}[Model-categorical transfer]
\label{rem:ainf2-heal-72}

Francis--Gaitsgory provide the enhanced Ran bar--cobar equivalence.
Proposition~\ref{prop:fg-ambient-properties} supplies the chosen model
presentation, and Theorem~\ref{thm:factorization-properad-transfer}
states the five additional hypotheses for factorization-enriched
properads.  The presentation clause $H_{\mathrm{HR}}$
(Theorem~\ref{thm:hackney-robertson-model}, Conjectured) posits the
simplicial properad model structure appearing at the base of this
transfer; \cite{HRY17} constructs the graphical-set homotopy theory
that evidences it.
\end{remark}

\begin{remark}[Chain realization]
\label{rem:ainf2-heal-73}

The enhanced Ran equivalence is constructed from the universal
properties of enhanced bar and cobar.  A strict chain realization uses
either a specified split contraction in the chosen model enhancement
or a Cousin resolution compatible with the diagonal stratification.
On the quadratic locus, the acyclic twisted complex in
$H_{\mathrm{CL}}(A,A^{\mathrm i},\tau_{\mathrm i})$ supplies the corresponding
contraction.
\end{remark}

\begin{remark}[Truncation-wise Hochschild finiteness]
\label{rem:ainf2-heal-74}

The convergence package uses truncation-wise finiteness with
conformal-weight control.  For a chiral algebra with positive-energy conformal
weight grading
$\cA = \bigoplus_{n\geq 0} \cA_n$ and finite-dimensional
weight spaces $\cA_n$, the bar complex $\bar B(\cA)$ is
truncation-wise finite (each conformal weight stratum
$\bar B(\cA)_{\leq N}$ has finite-dimensional cohomology), and
Theorem~\ref{thm:A-infinity-2} applies to the pro-nilpotent completion
$\hat\cA := \varprojlim_N \cA / \cA_{\geq N}$. This weakening captures
Virasoro (weight space $\cA_n = S^n \Vir_0 / \mathrm{null}$ with
finite dimension), Yangian (weight grading by spectral-parameter
order), and critical KM (weight grading by conformal weight of
the current algebra).  The resulting equivalence lies in the
weight-completed ambient specified by $H_{\mathrm{aug}}$ and
$H_{\mathrm{ML}}$.
\end{remark}

\begin{remark}[Correct cross-volume citation]
\label{rem:ainf2-heal-195}

The precise cross-volume references are
Theorem~\ref{thm:koszul-reflection}\ref{KR-i} for the published Ran
equivalence, Theorem~\ref{thm:A-infinity-2} for its conditional
factorization restriction, Theorem~\ref{thm:A-infinity-2}\ref{A2-i}
for the $H_{\mathrm{prop}}$ properadic envelope, and
Corollary~\ref{cor:classical-A-from-A-infinity-2} for the chain model
at a marked point.
\end{remark}

%================================================================
\section{Formal smoothness of chirally Koszul algebras}\label{sec:ainf2-formal-smoothness}
%================================================================

\begin{corollary}[Cotangent criterion for properadic-envelope formal
smoothness; \ClaimStatusConditional]
\label{cor:chiral-KK-formal-smoothness}
\index{formal smoothness!chiral Koszul implies}
\index{Koszul Kontsevich!formal smoothness}
\textup{[Type signature: Open quadrant; factorization-properad
presentation; Beilinson level $1\leftrightarrow2$ deformation lane;
hypothesis package $H_{\mathrm{prop}}$ plus cotangent Ext-vanishing.]}
Let $\cA$ be a chirally Koszul factorization algebra on $X$
\textup{(}Definition~\textup{\ref{def:chiral-koszul-pair}}), and assume
$H_{\mathrm{prop}}$.  Let $L_{\cA}^{\mathrm{prop}}$ be its cotangent
complex in
$\FactProp^{\aug}_{\langle\chirAss\rangle}(X)$.  If
\[
 \operatorname{Ext}^1_{\cA}(L_{\cA}^{\mathrm{prop}},M)=0
\]
for every coefficient module $M$ arising from a square-zero
extension, then $\cA$ is formally smooth in this properadic envelope.
For a square-zero extension $R\to R_0$ with kernel $M$ and an
$R_0$-valued point through $\cA$, the space of lifts is a torsor for
$\operatorname{Map}_{\cA}(L_{\cA}^{\mathrm{prop}},M)$; it is
contractible when this derived mapping space is contractible.
\end{corollary}

\begin{proof}
Cotangent deformation theory places the obstruction to a lift in
$\operatorname{Ext}^1_{\cA}(L_{\cA}^{\mathrm{prop}},M)$.  Its
vanishing gives a lift.  The derived mapping space from the cotangent
complex to $M$ acts simply transitively on the lift space, which gives
the asserted torsor and contractibility statements.  The package
$H_{\mathrm{prop}}$ supplies the model category in which this
cotangent complex and its derived mapping space are formed.
\end{proof}

\begin{remark}[Compatibility with the five-theorem architecture]
\label{rem:ainf2-architecture}
The unified Theorem~A
(Theorem~\ref{thm:koszul-reflection}) is the bar--cobar spine of the
volume. It occupies the first place in the five-package architecture
$A,B,C,D,H$: Theorems~B and~D use its ambient-qualified inversion and
obstruction theory, Theorem~C uses its Koszul-dual Verdier surface,
and Theorem~H uses its chiral Hochschild control on the Koszul/generic
locus.  Corollary~\ref{cor:ainf2-downstream-list} records the five
typed outputs: Ran reconstruction, factorization closure, completion,
properadic transfer, and ordered descent.
\end{remark}

%================================================================
\section{Platonic obstructions: four named conjectures}\label{sec:ainf2-platonic-obstructions}
%================================================================

Four conjectures govern the extension of Koszul reflection to the full
category $\Alg^{\fact,\aug}_X$.

\begin{conjecture}[$\Pi 1$: Francis--Gaitsgory transfer]
\label{conj:fg-transfer}
\ClaimStatusConjectured
The $(\infty,1)$-categorical bar--cobar Quillen equivalence transfers
from $\IndCoh(\Ran(X))$ to $\Alg^{\fact,\aug,\comp}_X$ and its
coalgebra dual through the factorization comparison functor.
\end{conjecture}
Lurie supplies the topological transfer theorem
\textup{(}Lurie HA~\S5.5\textup{)}. This
theorem-level use in the chiral setting requires a named transfer
proposition \texttt{prop:chiral-quillen-transfer}.

\begin{conjecture}[$\Pi 2$: $E_n$-bar--cobar at higher chiral
dimension]
\label{conj:en-bar-cobar-higher-d}
\ClaimStatusConjectured
For $d\geq 2$, the $E_n$-bar functor on $n$-dimensional factorization
$E_n$-algebras admits a left adjoint (cobar) that is a Quillen
equivalence onto conilpotent CDG-coalgebras over the appropriate
$n$-dimensional Maurer--Cartan complex. The $n=1$ case is
Theorem~\ref{thm:koszul-reflection}.
\end{conjecture}
This is the load-bearing direction connecting Vol~I to Vol~III's
CY-to-chiral functor. Currently open in the chiral setting (proved in
Lurie~HA~\S5.5.3 for the topological setting).

\begin{conjecture}[$\Pi 3$: Lagrangian--Koszul converse]
\label{conj:lagrangian-koszul-converse}
\ClaimStatusConjectured
Lagrangian transversality of $M_{\mathrm{comp}}$ at $[\cA]$ implies
chiral Koszulness of $\cA$.
\end{conjecture}
This requires the chiral PTVV shifted-symplectic construction.

\begin{conjecture}[$\Pi 4$: Unbounded-rank Koszul reflection]
\label{conj:unbounded-rank-koszul-reflection}
\ClaimStatusConjectured
On the locus of complete augmented chiral algebras with
arbitrary graded bar pieces, the Koszul reflection of
Theorem~\ref{thm:koszul-reflection} holds in the coderived category.
Its chain-level realization is the locus on which the pro-coradical
tower satisfies $H_{\mathrm{ML}}$.
\end{conjecture}
Finite-dimensional weight pieces give Mittag--Leffler towers for the
standard Heisenberg, Kac--Moody, Virasoro, and $W_N$ presentations.
The conjecture identifies $H_{\mathrm{ML}}$ as the chain-level frontier
inside the full completed category.

%================================================================
\section{Comparison with independent constructions}\label{sec:ainf2-hziv}
%================================================================

\begin{remark}[Sources for the five typed outputs]
\label{rem:ainf2-hziv}
\index{Theorem A-infinity-2!comparison constructions}
The five outputs of Corollary~\ref{cor:ainf2-downstream-list} have the
following sources:
\begin{itemize}
\item Francis--Gaitsgory's Proposition~4.1.2, together with the
pro-nilpotence argument in the proof of their Theorem~5.1.1,
supplies the enhanced Ran equivalence \cite{Francis2012}.
\item $H_{\mathrm{fact}}$ supplies the factorization closure and the
Fulton--MacPherson nearby-cycle extension.
\item Loday--Vallette \cite[Corollary~2.3.4]{LV12} supply the ordinary
chain-level unit and counit equivalences; $H_{\mathrm{conv}}$ carries
them through completion.
\item The presentation clause $H_{\mathrm{HR}}$
(Theorem~\ref{thm:hackney-robertson-model}, Conjectured) posits the
cofibrantly generated model structure on simplicial properads; the
graphical-set homotopy theory of \cite{HRY17} is its evidence.
\item $H_{\mathrm{prop}}$ supplies the factorization-enriched transfer
and the connected-graph preservation of the derived unit and counit.
\end{itemize}
\end{remark}

\begin{remark}[Independent constructions for bar-cobar reconstruction]
\label{rem:ainf2-hziv-koszul-reflection}
\index{bar-cobar reconstruction!comparison constructions}
The reconstruction statement has two published realizations and two
conditional extensions:
\begin{itemize}
\item Francis--Gaitsgory give enhanced reconstruction in the
pro-nilpotent Ran ambient.
\item Loday--Vallette give the ordinary associative chain models
$T^c(s^{-1}\bar A)$ and $T(s\bar C)$.
\item $H_{\mathrm{fact}}+H_{\mathrm{conv}}$ give the factorization and
completed realizations.
\item $H_{\mathrm{prop}}$ gives the properadic-envelope realization.
\end{itemize}
\end{remark}

%================================================================
\section*{Theorem~A as the algebraic structure on the KZ functor}
%================================================================

One adjunction, one reconstruction counit, one Verdier-dualized
duality functor.  Theorem~\ref{thm:A-infinity-2} acts by
factorization restriction, fiber realization, and pro-conilpotent
completion.  Verdier duality is the $K_X$ branch; a quantum $R$-matrix
is the ordered-descent input; the equations
$[\alpha_{\mathrm{BC}}]=[\beta_{\mathrm{nodal}}]=0$ give the
Deligne--Mumford family extension.  Theorem~A enters the
five-theorem spine of Volume~I as the algebraic structure on the KZ
functor of Chapter~\ref{chap:climax-platonic}: the pullback identity
$\bar\partial = \KZ^{*}(\nabla_{\Arnold})$ is the bar differential read
through the Koszul reflection.

\section{K3 Hall--Drinfeld input and Theorem~A scope}
\label{sec:tainf2-supplement}

\begin{remark}[$A_\infty$-structure on $\mathbf H_{\Delta_5}$ as external input]
\label{rem:tainf2-1}
Theorem~A supplies the completed bar--cobar ambient for a fixed chiral
shadow of a supplied factorization algebra.  The K3/Hall--Drinfeld
input consists of the object $\mathbf H_{\Delta_5}$, its comparison
with a Drinfeld Yangian, and its higher products.  Suppose Volumes
II--III supply a completed K3 chiral shadow
$\mathbf H_{\Delta_5}$ with $A_\infty$ structure maps
\[
\mu_n: \mathbf H_{\Delta_5}^{\otimes n}\to \mathbf H_{\Delta_5},\qquad n\ge 2,
\]
and verify $H_{\mathrm{fact}}+H_{\mathrm{conv}}$ together with the
relevant $R$-descent hypotheses.  The proposed first products
\[
\mu_2=\text{Hall multiplication in }\mathrm{CoHA}_{K3\times E},\qquad
\mu_3=\text{Drinfeld-associator coboundary at order }\hbar,
\]
and
\[
\mu_4=\text{Igusa-form twist at order }\hbar^2
\quad\text{from}\quad \Phi_{10}/\eta^{24}
\]
belong to the K3/Hall--Drinfeld input surface. The order
$\hbar^{\leq3}$ pentagon coboundary
\[
\phi^{(3)} = \zeta(3)c_{\mathrm{symm}}
+ \tfrac{25}{3}c_{\mathrm{timelike}}
+ (\Phi_{10}/\eta^{24})c_{\Phi_{10}}
\]
is likewise external structure data.  The completed bar--cobar
adjunction transports this supplied datum.
\end{remark}
